\subsection{Integral vanishing, cyclotomic descent and the number of exceptions}
\label{newcoh:section}

For a number field \(F\), let \(D_v,I_v\) denote absolute decomposition
and inertia groups at finite places in a fixed algebraic closure. For a
discrete \(\ell\)-primary \(G_F\)-module \(A\), write \(S_\ell(F,A)\)
for the classes in continuous \(H^1(F,A)\) which restrict to zero on
every \(I_v\) and on every \(D_v\) above \(\ell\). By conjugacy the
definition is independent of the choices of prolongations. When \(A\)
is unramified away from \(\ell\), this is exactly the kernel on
\(H^1(G_{F,S},A)\) used in Section~\ref{sup:length}. Additional
decomposition conditions at the primes of \(N\) define the subgroup
used at level \(\ell N\).
\src{src/Development/AutomorphicArithmetic/AutomorphicVarieties/CharacterLevelBounds/ShaFamMono.lean}{17}{coefficient action and the two levels}.

\begin{proposition}[The trivial-character integral vanishing]\label{newcoh:zero}
Suppose \(\ell\ge5\) is prime, \(e\ge2\) is even and \(\ell-1\mid 2e\).
Let \(O\) be the ring of integers in a finite extension of
\(\mathbf Q_\ell\). For \(A=(O\otimes_{\mathbf Z_\ell}
\mathbf Q_\ell/\mathbf Z_\ell)(-e)\) one has \(S_\ell(\mathbf Q,A)=0\).
In particular its subgroup at level \(\ell N\) has length zero for
every \(N\).
\end{proposition}
\begin{proof}
Every continuous cocycle on an absolute Galois group takes only finitely
many values in \(A\), by compactness. Its class is therefore killed by
some power of \(\ell\). It suffices to show that a class \(x\) killed
by \(\ell\) is zero. Fix a place above \(\ell\). Subtract a coboundary
to obtain a representative \(z\) vanishing on \(D_\ell\). Write
\(\ell z(g)=g b-b\). Then \(b\in A^{D_\ell}\). Every extension
\(\mathbf Q_\ell(\mu_{\ell^n})/\mathbf Q_\ell\) is totally ramified
of degree \(\varphi(\ell^n)\), as follows from the Eisenstein
polynomial \(\Phi_{\ell^n}(1+X)\). Hence the local cyclotomic image is
the full global image, so \(b\in A^{G_{\mathbf Q}}\). It follows that
\(z\) actually has values in \(A[\ell]\).

On \(A[\ell]\) the action is via the scalar character
\(\bar\chi_\ell^{-e}\). Its order divides two. If it is trivial put
\(K=\mathbf Q\). Otherwise \(e\) is an odd multiple of
\((\ell-1)/2\). Since \(e\) is even, \(\ell\equiv1\bmod4\), and the
quadratic Gauss sum identifies the fixed field of its kernel as
\(K=\mathbf Q(\sqrt{\ell})\), of discriminant \(\ell\).
Now \(z|_{G_K}\) is a homomorphism. It vanishes on inertia away from
\(\ell\), where the action on all of \(A\) is trivial and the
cohomology restriction of \(x\) is zero. At \(\ell\) it vanishes on
\(D_\ell\cap G_K\) by normalization. For its conjugates use
\[
 z(hgh^{-1})=h z(g)+(1-hgh^{-1})z(h).
\]
The second term is zero for \(g\in G_K\), and the first term is zero
when \(z(g)=0\). Real inertia has order two and cannot contribute.
Global class field theory now makes this homomorphism factor through
the ordinary ideal class group of \(K\).

For \(K=\mathbf Q\) this group is trivial. For the quadratic field put
\(y=\sqrt{\ell}\). Minkowski's bound supplies each class with an
integral ideal of norm at most \(y\). There are at most \(\tau(n)\)
ideals of norm \(n\), by the three quadratic splitting possibilities
and unique ideal factorization. Hence
\[
 h_K\le\sum_{n\le y}\tau(n)=\sum_{a\le y}\lfloor y/a\rfloor
 \le y(1+\log y)<y^2=\ell .
\]
Thus \(\ell\nmid h_K\), and the restriction vanishes in both cases.
The cocycle factors through \(\operatorname{Gal}(K/\mathbf Q)\) with
values in \(A[\ell]\). This quotient has order dividing two, which
acts invertibly on \(A[\ell]\), so its first cohomology vanishes by
averaging. This proves \(x=0\).
\end{proof}
The local normalization is what makes reduction to \(A[\ell]\)
legitimate even in the residually trivial case. This gives a direct
proof of the source's \(\mathrm{PropZ}\).
\src{src/Development/AutomorphicArithmetic/AutomorphicVarieties/LevelCohomology/PropZ.lean}{11}{two residual characters},
\src{src/Development/AutomorphicArithmetic/AutomorphicVarieties/LevelCohomology/HQuad.lean}{11}{quadratic class bound}.

\paragraph{The split unramified module.}
Assume \(F\) contains \(\mu_\ell\) for odd \(\ell\) and \(\mu_4\) for
\(\ell=2\). Put \(F_\infty=F(\mu_{\ell^\infty})\),
\(\Gamma=\operatorname{Gal}(F_\infty/F)\) and \([F_m:F]=\ell^m\) in
this tower. Let \(X\) be the Galois group of the maximal abelian
pro-\(\ell\) extension of \(F_\infty\) unramified at every finite
place in which every place above \(\ell\) splits completely. Equivalently
it is the maximal abelian pro-\(\ell\) quotient of \(G_{F_\infty}\)
which kills all finite inertia groups and all those decomposition
groups. This definition also specifies what unramified means over the
infinite field. Conjugation induces an action of \(\Gamma\).
Choose a generator \(\gamma\) of \(\Gamma\simeq\mathbf Z_\ell\), write
\(g\) for its action on \(X\), and put \(u=\chi_\ell(\gamma)\).
With \(c=v_\ell(u-1)\) one has \(c\ge1\) and \(c\ge2\) at \(\ell=2\).
The group \(\Gamma\) is an open subgroup of the principal units
under the root-of-unity hypotheses.

\begin{lemma}[Descent with the local conditions]\label{newcoh:desc}
For \(e\ge1\), the Pontryagin dual of
\(S_\ell(F_m,(\mathbf Q_\ell/\mathbf Z_\ell)(-e))\) is
\[
 X\big/\big((u^eg)^{\ell^m}-1\big)X .
\]
\end{lemma}
\begin{proof}
Put \(A_e=(\mathbf Q_\ell/\mathbf Z_\ell)(-e)\).
The quotient \(\Gamma_m=\operatorname{Gal}(F_\infty/F_m)\) is
procyclic with generator \(\gamma^{\ell^m}\). Its first cohomology on
\(A_e\) is zero since \(u^{-e\ell^m}-1\ne0\) acts surjectively on
\(\mathbf Q_\ell/\mathbf Z_\ell\). Its second cohomology is zero since
\(\mathbf Z_\ell\) has \(\ell\)-cohomological dimension one.
Inflation--restriction identifies \(H^1(F_m,A_e)\) with the
\(\Gamma_m\)-invariant continuous homomorphisms from \(G_{F_\infty}\)
to \(A_e\), on which \(G_{F_\infty}\) acts trivially.
The tower is unramified away from \(\ell\), so the inertia subgroups
there are unchanged. At \(w\mid\ell\) the local cyclotomic image is an
open subgroup of \(\Gamma_m\), since \((F_m)_w/\mathbf Q_\ell\)
is finite. The same procyclic first-cohomology calculation gives
injectivity of local restriction. Thus the finite-level local
conditions are exactly vanishing on the defining subgroups of \(X\).
We obtain \(\operatorname{Hom}_{\rm cont}(X,\mathbf Q_\ell/
\mathbf Z_\ell)(-e)^{\Gamma_m}\). Duality turns invariants into
coinvariants, and the dual twist acts by \(u^e\). The image of the
continuous endomorphism of the compact group \(X\) is closed, so its
ordinary quotient gives those coinvariants.
\end{proof}
\noindent
\src{src/Development/AutomorphicArithmetic/AutomorphicVarieties/GaloisLibrary/KummerTheory/Mk0IdMapCongr.lean}{90}{local conditions over the tower},
\src{src/Development/AutomorphicArithmetic/AutomorphicVarieties/NumberFieldTowers/ExistsDescends.lean}{26}{descent with the positive dual twist}.

\begin{proposition}[Integral growth and centres]\label{newcoh:centres}
Suppose \(X\) has at most \(R\ge1\) generators over \(\mathbf Z_\ell\)
and \(g\) is unipotent on \(X/\ell X\). Set
\(j=\lfloor\log_\ell R\rfloor+3\) and
\(L_m(e)=\operatorname{length}_{\mathbf Z_\ell}
\big(X/((u^eg)^{\ell^m}-1)X\big)\), allowing infinite length.
There are \(R\) elements \(s_1,\ldots,s_R\in\mathbf Z_\ell\),
independent of \(e,m\), such that
\begin{align}
 L_m(e)&\le L_j(e)+R(m-j)_+, \label{newcoh:growth}\\
 L_m(e)&\le\sum_{i=1}^R
 \left(\max(m,j)+c+v_\ell(e-s_i)\right). \label{newcoh:centreformula}
\end{align}
For each integer \(Z\ge1\), at most \(R\) values of \(1\le e\le Z\)
need be removed to bound every \(v_\ell(e-s_i)\) by
\(\lfloor\log_\ell Z\rfloor\) for the remaining \(e\).
\end{proposition}
\begin{proof}
Write \(h=u^eg\) and \(k=\lfloor\log_\ell R\rfloor+1\).
On \(X/\ell X\) the dimension is at most \(R\) and
\((h-1)^R=0\). In characteristic \(\ell\),
\(h^{\ell^k}-1=(h-1)^{\ell^k}=0\) there. Thus
\((h^{\ell^k}-1)X\subseteq\ell X\). Raising this endomorphism
plus the identity to the \(\ell\)-th power shows
\((h^{\ell^{k+1}}-1)X\subseteq\ell^2X\) by binomial expansion.
This inclusion persists at higher layers.

For \(w=h^{\ell^n}\), \(n\ge j\), put \(P=1+w+\cdots+w^{\ell-1}\).
Then \(Px\equiv\ell x\bmod\ell^2X\), whence \(PX\subseteq\ell X\)
and \(PX+\ell^2X=\ell X\). Nakayama gives \(PX=\ell X\).
Consequently
\[
 (w^\ell-1)X=\ell(w-1)X,\qquad
 (w-1)X/(w^\ell-1)X=(w-1)X/\ell(w-1)X .
\]
The last module has dimension at most \(R\), since \((w-1)X\) is a
quotient of \(X\). This proves the claimed one-step loss even when
\(X\) has \(\ell\)-torsion. Addition of lengths proves
\eqref{newcoh:growth} for \(m\ge j\). For \(m<j\) use the
divisibility \(T^{\ell^m}-1\mid T^{\ell^j}-1\).

Choose a surjection \(\mathbf Z_\ell^R\to X\) and lift the images of
the generators under \(g^{\ell^j}\) to columns of an integral matrix
\(C\). Put \(a=u^{\ell^j}\). Then \(X/((u^eg)^{\ell^j}-1)X\)
is a quotient of \(\operatorname{coker}(a^eC-I)\). By Smith normal form
the length of the latter equals \(v_\ell(\det(a^eC-I))\), with
\(v_\ell(0)=+\infty\). Let \(z_i\) be the eigenvalues of \(C\)
in a splitting field, counted with multiplicity. For a unit
eigenvalue \(z\), choose \(s\in\mathbf Z_\ell\) minimizing
\(\lvert a^s z-1\rvert_\ell\). Such \(s\) exists by compactness and
the continuous extension of powers on the principal units. For
integers \(e\ge0\), the ultrametric inequality gives
\[
 |a^e-a^s|_\ell\le\max(|a^e z-1|_\ell,|a^s z-1|_\ell)
 = |a^e z-1|_\ell .
\]
The binomial valuation formula from Section~\ref{sup:length}, extended
by continuity from integers, gives
\(v_\ell(a^e-a^s)=j+c+v_\ell(e-s)\), including equality at infinity.
For a nonunit eigenvalue \(|z|_\ell<1\) the valuation of \(a^e z-1\)
is zero and any centre suffices. All eigenvalues are integral.
Factoring the determinant proves the sum bound at \(j\), and
\eqref{newcoh:growth} proves \eqref{newcoh:centreformula}. Finally
integers with \(v_\ell(e-s_i)>\lfloor\log_\ell Z\rfloor\) form at
most one residue class modulo \(\ell^{1+\lfloor\log_\ell Z\rfloor}>Z\);
this class meets \([1,Z]\) at most once.
\end{proof}
In Lemma~\ref{newcoh:desc} the action on the finite vector space
\(X/\ell X\) factors through a finite \(\ell\)-group, which gives
the required unipotence. The centres are attached to a chosen base
field. Only their number must be uniform as the auxiliary primes vary.
\src{src/Development/AutomorphicArithmetic/AutomorphicVarieties/AdditiveValuationNorms/ExistsCentresLevel.lean}{19}{integral tower inequalities},
\src{src/Development/AutomorphicArithmetic/AutomorphicVarieties/LayeredIsomorphisms/Exc.lean}{8}{exceptional set from centres}.

\paragraph{Uniformity for the moving quadratic fields.}
Let \(M/\mathbf Q\) be the fixed abelian target field, \(d_M=[M:\mathbf Q]\),
and let \(q\equiv3\bmod8\), \(d\equiv7\bmod16\) be distinct primes
coprime to \(2\) and to the chosen modulus of \(M\).
Put \(B=M(\mu_{2^\infty})\). The unramified elementary
two-characters of \(G_B\) have dimension at most a constant \(g_M\)
depending only on \(M\). Indeed, Ferrero--Washington gives
\(\mu=0\) for the fixed abelian base \(M(i)\). The cyclotomic
Iwasawa structure theorem then gives a bound for the two-ranks of the
class groups throughout this fixed tower. Class field theory gives
the same bound on unramified two-characters in its finite layers.
For the passage to the union, every finite collection of unramified
characters at infinite level descends unramified to one finite level.
To see this for one character, extend its locally constant function on
the closed subgroup \(G_B\) to a function factoring through a finite
quotient of \(G_{M(i)}\). On some finite layer its image subgroup in
this finite quotient is already that of \(G_B\), so the extension
there is a homomorphism. Only finitely many places of \(M(i)\) can
ramify in that quotient. For each of them choose one absolute inertia
group and its finitely many conjugates in the quotient. More precisely,
apply stabilization inside each conjugate of the absolute inertia
group to the descending compact intersections with the layer groups.
Their images eventually equal the image of the intersection with
\(G_B\). One representative per coset of the open normal kernel of
the finite quotient suffices, since the finite layer groups are normal
over \(M(i)\) and conjugation by that kernel acts trivially in the
quotient. The extension is therefore unramified on a sufficiently
high layer. Taking the maximum for a finite collection proves the
bound at infinite level.
See Ferrero--Washington, Ann. of Math. 109 (1979)~\cite{ferreroWashington}.
\src{src/Development/AutomorphicArithmetic/AutomorphicVarieties/PublishedTheorems/FerreroWashington.lean}{7}{finite-layer rank formulation},
\src{src/Development/AutomorphicArithmetic/AutomorphicVarieties/QuadraticExtensions/LevUAnti.lean}{9}{passage to the union}.

We compute explicitly how much this dimension can increase when
adjoining \(\sqrt{-q},\sqrt{-d}\). In the full two-power cyclotomic
tower over \(\mathbf Q\), the number of primes over an odd prime \(p\)
is the index of the closed subgroup generated by \(p\) in
\(\mathbf Z_2^\times\). For the two congruence classes at hand
\(v_2(q^2-1)=3\), \(v_2(d^2-1)=4\). The even powers of each
prime therefore generate \(1+2^3\mathbf Z_2\) and
\(1+2^4\mathbf Z_2\), respectively, by the binomial valuation formula.
Each prime itself is \(3\bmod4\), so adjoining its odd powers gives
indices \(2\) and \(4\) in \(\mathbf Z_2^\times\). Base change to
\(B\) multiplies the number of places by at most \(d_M\).
Allowing ramification at \(q,d\) thus raises the dimension of
two-characters by at most \(6d_M\). Indeed a character's restriction
to each new inertia group is specified by one bit, since wild inertia
has odd residue characteristic and the tame pro-two quotient is
procyclic. The kernel of these evaluations is unramified everywhere.

Each of the two quadratic adjunctions multiplies the resulting dimension
by at most two. To see this, quotient the absolute group by the closed
normal subgroup generated by inertia outside \(q,d\). Each adjunction
is unramified there, so defines a subgroup \(U\) of index at most two
in the resulting group \(V\). Normality of \(U\) and the use of all
conjugate inertia groups ensure that the same quotient inside \(U\)
imposes all the required inertia conditions. For index two let
\(W=H^1(U,\mathbf F_2)\). The inflation--restriction sequence gives
\(\dim W^{V/U}\le\dim H^1(V,\mathbf F_2)\), since the inflation kernel
in degree one contributes one dimension and the transgression has
target \(H^2(C_2,\mathbf F_2)\) of dimension one. On \(W\) the
operator \(\sigma-1\) for the nonidentity coset satisfies
\((\sigma-1)^2=0\), so \(\dim W\le2\dim W^{V/U}\).
It follows that the split unramified \(X\) for
\(M(i,\sqrt{-q},\sqrt{-d})\) has generator bound
\[
 R_2=4(g_M+6d_M).
\]
Indeed \((X/2X)^\vee\) is a subspace of the characters just counted,
and lifts of any basis of \(X/2X\) generate the compact abelian
pro-two group \(X\) over \(\mathbf Z_2\). The closed span of the
lifts has zero quotient by compact Nakayama; it equals their algebraic
\(\mathbf Z_2\)-span by compactness. The same fixed-field argument
before the quadratic adjunctions applies at each of the finitely many
fixed odd primes with its own abelian base containing \(\mu_\ell\).
Consequently the cardinality bound used to choose the window precedes
the choice of \(q,d\).
\src{src/Development/AutomorphicArithmetic/AutomorphicVarieties/QuadraticExtensions/IsAlgClosureRat.lean}{69}{moving two-rank and its exact constant},
\src{src/Development/AutomorphicArithmetic/AutomorphicVarieties/LayeredIsomorphisms/ThmBabOf.lean}{27}{sum of the uniform ranks}.

\subsection{An independent bound for the low layer, including the small primes}
\label{newcoh:ktheory}

\begin{theorem}[Uniform order of the low-layer group]\label{newcoh:lemmaK}
Let \(F\) be totally imaginary of degree \(n\), with absolute
discriminant \(D\). Let \(E\ge1\) and let \(\ell\) satisfy the
root-of-unity hypotheses of Section~\ref{sup:length}. With \(C=32\)
one has
\[
 H^2(G_{F,S_\ell},\mathbf Z_\ell(E+1))\ \text{is finite},\qquad
 \log\#H^2(G_{F,S_\ell},\mathbf Z_\ell(E+1))
 \le (2D)^{(2En)^C}.
\]
\end{theorem}
\begin{proof}
Put \(m=2E\), \(N=2m+1\ge5\), \(U=\operatorname{Spec}
\mathcal O_F[1/\ell]\), and \(A=K_m(\mathcal O_F)\{\ell\}\), where
braces indicate the primary subgroup. Quillen's finite generation and
Borel's rational rank formula imply that \(K_m(\mathcal O_F)\) is
finite. The Quillen--Lichtenbaum comparison for number rings gives
\[
 A\simeq H^2_{\mathrm{\acute et}}(U,\mathbf Z_\ell(E+1)).
\]
Here \(E+1\ge2\). At the prime two all correction terms arising from
real embeddings vanish. Inverting \(\ell\) in the \(K\)-group also
gives the same primary group, since the finite residue fields above
\(\ell\) have vanishing \(K_{2E}\) and \(K_{2E-1}\) of order
\((\#k_v)^E-1\). These comparisons hold before taking orders.
The rank calculation is Borel's Proposition~12.2~\cite{borelStable}.
For the comparison see the norm-residue theorem~\cite{voevodskyNorm},
Kahn's Theorem~1 at two~\cite{kahnTwo}, and Weibel, Chapter~VI,
Theorems~8.2 and 8.4~\cite{weibelKbook}.

For the desired upper bound an injection of Galois cohomology into
this group suffices. For each finite locally constant coefficient sheaf
on \(U\), the fundamental-group comparison in degree two is injective.
Indeed degree one on the universal pro-finite-\'etale cover vanishes,
because each torsor for a finite constant abelian group is trivialized
on a finite \'etale cover. The five-term sequence of that cover gives
the injection, and \(\pi_1(U)=G_{F,S_\ell}\). Apply this to
\(\mu_{\ell^a}^{\otimes(E+1)}\) and take inverse limits. The finite
degree-one groups on the number ring and its fundamental group satisfy
the Mittag--Leffler condition, so no derived limit in degree one enters.
Their finiteness also follows by trivializing the finite action in a
finite extension and applying class field theory to extensions of
bounded degree unramified outside \(S_\ell\).

\emph{The number of generators at the small primes.}
If \(r=\dim_{\mathbf F_\ell}A/\ell A\), reduction of the displayed
\'etale cohomology injects \(A/\ell A\) into
\(H^2_{\mathrm{\acute et}}(U,\mu_\ell^{\otimes(E+1)})\).
Choosing a primitive \(\ell\)-th root over \(F\) identifies the latter
coefficient sheaf with \(\mu_\ell\). By the Kummer sequence,
\[
 0\longrightarrow\operatorname{Pic}(U)/\ell
 \longrightarrow H^2_{\mathrm{\acute et}}(U,\mu_\ell)
 \longrightarrow\operatorname{Br}(U)[\ell]\longrightarrow0 .
\]
The class group surjects onto \(\operatorname{Pic}(U)\). The local
invariant sequence for the Brauer group embeds
\(\operatorname{Br}(U)[\ell]\) in the \(\ell\)-torsion of the
sum-zero subgroup of \(\bigoplus_{v\mid\ell}\mathbf Q/\mathbf Z\).
All archimedean places are complex. In particular
\[
 r\log\ell\le\log h_F+(n-1)\log\ell .                 \tag{*}
\]
Minkowski supplies each ideal class with an integral ideal of norm at
most \(2^n\sqrt D\). For \(y\ge1\) the number of such ideals of norm
at most \(y\) is at most \(y^2\zeta_F(2)\le y^2\zeta(2)^n\), since
each rational prime has at most \(n\) Euler factors of residue degree
at least one. Thus \(\log h_F\le\log D+O(n)\). Also
\(\ell\le n+1\). The constants introduced here are absolute.

\emph{The integral Hurewicz map and transfer.}
The identity component \(B\mathrm{GL}(\mathcal O_F)^+\) of the
algebraic \(K\)-theory infinite loop space is connected. Apply Mathew's
unstable Hurewicz theorem to this component in degree \(m\), with
connectivity parameter one. Its \(\ell\)-primary kernel is killed by
\(\ell^{m+3}\). Indeed his two bounds on the exponent of \(\ell\) are
\(\lceil(m-1)/2\rceil+3\) at two and
\(\lceil(m+2)/(2\ell-2)\rceil+1\) at odd primes. The plus
construction preserves integral homology. By homological stability
for the ring of integers (a Dedekind domain of stable rank at most
two), \(H_m(\mathrm{GL}(\mathcal O_F),\mathbf Z)\) identifies with
\(H_m(\Gamma,\mathbf Z)\) for
\(\Gamma=\mathrm{GL}_N(\mathcal O_F)\) and \(N=2m+1\).
See Mathew, Theorem~6.3~\cite{mathewHurewicz}, van der Kallen,
Theorem~4.11~\cite{vanderKallen}, and Bayer-Fluckiger, Emery and
Houriet, \S3.1~\cite{bayerEmeryHouriet}.

Let \(\Gamma'\) be the kernel of reduction modulo \(3\) in \(\Gamma\)
and \(I=[\Gamma:\Gamma']\le3^{nN^2}\). The group \(\Gamma'\) is
torsion-free. In the action on the underlying integral lattice, a
nonidentity element in it has the form \(1+3^tB\) for \(t\ge1\) and
an integral matrix \(B\) nonzero modulo three. Raising to a prime
exponent \(p\ne3\) leaves the first nonzero
three-adic term at \(t\); cubing puts it exactly at \(t+1\), by the
binomial formula. In particular there can be no element of prime order.
Compose the Hurewicz map with transfer to
\(H_m(\Gamma',\mathbf Z)\). If \(x\in A\) is in the kernel of this
composite, the pushforward of its transfer gives \(Ih(x)=0\), so
\(\ell^{m+3+v_\ell(I)}x=0\). A subgroup of a finite abelian
\(\ell\)-group on \(r\) generators killed by \(\ell^a\) has at most
\(\ell^{ar}\) elements, by the elementary divisors. We have proved
\begin{equation}\label{newcoh:transfer}
 \log\#A\le
 (m+3+v_\ell(I))\,r\log\ell+
 \log\# H_m(\Gamma',\mathbf Z)_{\mathrm{tors}} .
\end{equation}
This bounds the transfer loss also for primes dividing \(I\).

\emph{Bounding the torsion of the congruence subgroup.}
Use Ash's contractible well-rounded retract of the space of positive
archimedean Hermitian forms on the free lattice \(\mathcal O_F^N\).
Normalize the minimum to one. Each cell has finite stabilizer and is
determined by the set of its minimal vectors \(v\), those with
\(Q(v)=1\). For a specified minimal set the stratum is the relative
interior of the convex face imposed by its equalities in the region
\(Q(w)\ge1\) for every nonzero lattice vector \(w\). Its boundary
adds minimal vectors. This polyhedral structure gives the subdivision
by flags of faces, ordered by reverse inclusion of their minimal
sets. Put
\(b=nN^2\), which bounds the dimension of the real vector space of
Hermitian forms. By Bayer-Fluckiger--Emery--Houriet there exists a
finite \(\Phi\subset\mathcal O_F^N\) into which every one of these
minimal sets can be moved by \(\Gamma\), with
\[
 |\Phi|\le N^{3N^2n^2} n^{5N^3n^4} D^{9N^3n^3}.
\]
Their result applies to free lattices, \(N\ge5\) and \(F\ne\mathbf Q\).
These all hold, including for \(E=1\).
See Bayer-Fluckiger--Emery--Houriet, Thm.~2.2 and Cor.~5.2~\cite{bayerEmeryHouriet}.

Here is a count for the barycentric subdivision which allows redundant
minimal-vector equations. In any flag move the minimal-vector set of
its smallest face into \(\Phi\). Every other face in the flag then
has its minimal-vector set in \(\Phi\). A single such set \(T\) is
encoded by at most \(b+1\) elements of \(T\). Choose a basis for the
span of the affine functions \(Q\mapsto Q(v)-1\), \(v\in T\), on
the \(b\)-dimensional ambient space (or its actual smaller space).
The equalities for this basis recover all minimal vectors at that
face. In fact for the affine solution space of these equations, each
vector in \(T\) has \(Q(v)=1\) identically, whereas for \(v\notin T\)
the chosen interior point of the face has \(Q(v)>1\) when \(v\ne0\).
Thus distinct faces cannot share the same encoding. A flag has at
most \(b+1\) members. Encoding each choice as a padded list from
\(\Phi\) with an extra termination symbol gives the bound
\[
 C_{\rm cell}=(|\Phi|+2)^{(b+2)^2}
\]
for the number of \(\Gamma\)-orbits of simplices in the subdivision.
Consequently the quotient by \(\Gamma'\) has at most \(I C_{\rm cell}\)
simplices in every degree. The free action on the contractible
complex gives a classifying space for \(\Gamma'\).

For the boundary in degree \(m+1\), each column has Euclidean norm
at most \(m+2\), since there are \(m+2\) faces counted with sign;
possible identifications in the quotient respect this upper bound.
The torsion order of the cokernel of an integral matrix of rank \(s\)
is the gcd of its full-rank minors. Hadamard bounds every such minor
by the product of \(s\) column norms. The torsion in homology embeds
in the torsion of this cokernel, so
\[
 \log\#H_m(\Gamma',\mathbf Z)_{\mathrm{tors}}
 \le I C_{\rm cell}\log(m+2).
\]
Finally put \(P=2En\ge4\), \(Y=\log(2D)\ge\log2\).
Then \(N=4E+1\le2P\), \(n,m\le P\), \(b\le P^4\), and
\(\log I\le2P^4\). Using \(\log N\le2P,\ \log n\le P\),
the displayed bound for \(\Phi\) gives
\(\log|\Phi|\le24P^5+40P^8+72P^6Y\le P^{12}Y\).
Thus \(\log C_{\rm cell}\le(2P^4)^2(2P^{12}Y)\le P^{22}Y\).
By \((*)\), \(r\log\ell\le\log D+n\log(4\zeta(2))+n\log(n+1)
\le 6P^2Y\), while \(m+3+v_\ell(I)\le5P^4\).
Equation~\eqref{newcoh:transfer} now gives
\[
 \log\#A\le30P^6Y+\exp(2P^4+P^{22}Y)\log(m+2)
 \le\exp(P^{32}Y).
\]
For the last inequality each summand is at most
\(\exp(P^{24}Y)\), and \(2\exp(P^{24}Y)\le\exp(P^{32}Y)\)
throughout \(P\ge4\). This proves the estimate.
\end{proof}

Theorem~\ref{newcoh:lemmaK} applies directly over the field at the
fixed layer \(j\) in \eqref{newcoh:growth}. No enlargement of that
field or uniformity loss in a cyclotomic descent is needed in this
proof of the low-layer estimate. Poitou--Tate identifies the dual of
the strict degree-one kernel with the kernel of the degree-two
localization map, hence \(L_j(e)\le
\operatorname{length}H^2(G_{F_j,S_\ell},\mathbf Z_\ell(e+1))\).
Inserting the degree and discriminant bounds of Section~\ref{sup:length}
therefore gives the exceptional contribution calculated there. The
original proof reaches the same bound through the integral homology
cycles in \eqref{sup:lengthcount}; the proof here uses a different
set of classical theorems.

\paragraph{The primitive homology count in the original argument.}
One can also derive its factor \(q^{r_2}T_{2E+1}T_{2E}\) directly.
Let \(d=2E+1\), \(q=\ell^\nu\), and take the stable integral homology of
\(\mathrm{GL}(\mathcal O_{L'})\). A finite collection of mod-\(q\)
cycles of varying matrix sizes can be moved into one stable matrix
size, by stabilization followed by the natural map. Suppose their
joint evaluation \(\Phi\) is a homomorphism on stable mod-\(q\)
homology taking values in a \(\mathbf Z/q\)-module and satisfying
\(\Phi\psi_b=b\Phi\), where \(\psi_b\) denotes the \(b\)-fold
block-diagonal operation. These are precisely the boundary,
stabilization and direct-sum requirements on the primitive families.
Use integral stability in degrees \(d,d-1\) to write \(T_i\) for the
torsion order in either stable homology or matrix size \(2d+4\).

Consider \(P=\prod_{j=2}^{d}(\psi_\ell-\ell^j)\) on rational
homology in degree \(d\), induced by its integral counterpart. Rational
homology of the connected infinite loop space is the free graded
commutative algebra on its rational homotopy under Hurewicz, by the
rational Hurewicz theorem for connected homotopy-commutative
\(H\)-spaces. The product is block sum; its \(b\)-fold power operation
acts on a monomial of \(a\) positive-degree primitives by \(b^a\).
Since \(a\le d\), the image of \(P\) is contained in the primitive
subspace. For the totally imaginary field \(L'\), Borel's rank theorem
gives dimension \(r_2(L')\) in this odd degree \(d\ge3\).
Thus the integral image has rank at most \(r_2(L')\), and its torsion
has order at most \(T_d\). For any integral class \(y\),
\[
 \Phi(\overline{Py})=
 \ell^{d-1}\prod_{j=2}^{d}(1-\ell^{j-1})\,\Phi(\bar y).
\]
The product is a unit modulo \(q\). Hence the image of
\(\ell^{d-1}\Phi\) on reductions of integral classes has size at most
\(q^{r_2(L')}T_d\). The coefficient exact sequence identifies the
quotient of \(H_d(-,\mathbf Z/q)\) by integral reductions with
\(H_{d-1}(-,\mathbf Z)[q]\), of order at most \(T_{d-1}\).
Evaluating on representatives for this quotient multiplies the number
of possibilities by at most \(T_{d-1}\). Thus on \emph{all} mod-\(q\) cycles
\[
 \#\operatorname{im}\bigl(\ell^{d-1}\Phi\bigr)
 \le q^{r_2(L')}T_dT_{d-1}.
\]
This includes cycles with nonzero Bockstein and applies to the joint
product of residue evaluations used in \eqref{sup:lengthcount}. The
scalar is essential when \(q=\ell\) and \(d\ge2\), for then every
scaled evaluation is zero regardless of the unscaled image.
See Borel, \S\S10.6 and 12.1--12.2~\cite{borelStable}.
\src{src/Development/AutomorphicArithmetic/AutomorphicVarieties/ModularityArgument/ClassCounting/SUpperOf.lean}{41}{scaled map on all cycles and count by Bockstein fibres}.

The residual power test in that argument has a short field-theoretic
proof. Let \(F/L\) be finite, \(\mu_q\subset F\), and
\(\beta\in F^\times\). Suppose \(\beta\) is a \(q\)-th power in the
residue fields at all but finitely many primes of \(F\) of relative
residue degree one over \(L\). Then \(\beta\in F^{\times q}\).
Indeed, if not, \(F(\beta^{1/q})/F\) is a nontrivial cyclic Kummer
extension. Take the normal closure over \(L\) and an automorphism
fixing \(F\) which induces a nonidentity element of the Kummer
group. Chebotarev gives primes of \(L\) outside the exclusions and
prolongations with exactly this Frobenius. The primes of \(F\) below
these prolongations have residue degree one, since Frobenius fixes
\(F\) pointwise. It multiplies any \(q\)-th root of \(\beta\) by the
same nontrivial root of unity. Away from the primes of \(q\), \(\beta\)
and its denominator, these roots remain distinct on reduction, so
the residue polynomial has no root in the residue field of \(F\).
This contradiction proves the assertion.
\src{src/Development/EtaleCohomology/CyclotomicDiscriminants/SDetectCut.lean}{101}{residual power criterion}.

The original route additionally constructs cycles whose evaluations
encode twisted norms of actual divisors; the abstract evaluation-count
inequality just proved requires such an encoding to bound class-group
characters. Theorem~\ref{newcoh:lemmaK} yields the required cohomological
conclusion independently of that construction. The higher-rank
automorphic estimates enter elsewhere. For example, when \(N\) is
the host modulus, \(\psi(-1)=(-1)^e\), \(2\le e<t\), and the
prime-to-\(\ell\) part of \(\psi\) after deleting its local
\(\ell\)-component has even order, case II applies if either
\(\ell\) lies in the host's fixed exceptional prime set or
\(\ell-1\mid2e\). At odd \(\ell\), for the admitted hosts the
required estimate is finiteness of the strict length at level
\(\ell N\) together with
\[
 \frac{\log\ell}{e(O/\mathbf Z_\ell)}
 \operatorname{length}_O S_{\ell N}
   \big(\mathbf Q,(\operatorname{Frac}O/O)(\psi^{-1}\chi_\ell^{-e})\big)
 \le (\ell N e)^{C_\chi}.
\]
Here \(S_{\ell N}\) imposes zero on all inertia and on \(D_v\)
for \(v\mid\ell N\), \(\psi\) occurs in the host character set,
and \(C_\chi\) depends only on the fixed target. Case III requires the
same inequality away from the exceptional prime set when
\(\ell-1\mid2e\) and \(\psi\ne1\) has odd order, for the applicable
positive twists. The cited proof uses the adjoint lifting theorem,
the symplectic rank-four counting result, and the residual Bernoulli
test. The low-layer theorem has a different dependence on \(e\) and
\(N\). Propositions~\ref{famsel:hostinertia} and~\ref{famsel:splitray}
prove vanishing in further cases at every congruence depth. The
complementary budget in Section~\ref{famsel:complement} uses the
cited polynomial bounds on the remaining entries.
\src{src/Development/AutomorphicArithmetic/AutomorphicVarieties/HeightAndVanishing/Counts.lean}{60}{quantified normalized bounds},
\src{src/Development/AutomorphicArithmetic/AutomorphicVarieties/ModularityArgument/ThresholdArguments/DenseThresholdCases/AdjointLiftingOfCohomology.lean}{33}{automorphic and residual inputs for cases II and III}.

There is a parity obstruction to using the ordinary special value as
the nonzero integer in this comparison. For either primitive inducer
\(\varphi\) of \(\psi\) or \(\psi^{-1}\), let \(f\) be its conductor
and \(m=e+1\). If \(f>1\), in the expression
\[
 B_{m,\varphi}=f^{m-1}\sum_{a=1}^{f}\varphi(a)B_m(a/f)
\]
the substitution \(a\mapsto f-a\) on units multiplies the sum by
\(\varphi(-1)(-1)^m=-1\), since \(B_m(1-x)=(-1)^m B_m(x)\).
Thus \(B_{m,\varphi}=0\). At \(f=1\), \(e\) is even and positive,
so \(m\ge3\) is odd and \(B_m(1)=B_m(0)=0\), by the same reflection
and \(B_m(1)-B_m(0)=0\) for \(m>1\). Hence \(L(-e,\varphi)=0\)
in these branches. Valuations of this value supply no finite bound
for the strict lengths. The trace products below require an actual
nonzero integer.

\begin{proposition}[Finiteness and prime support before the effective estimate]
\label{newcoh:fullsupport}
Fix \(N,e\ge1\), let \(F=\mathbf Q(\mu_{4N})\), \(n=[F:\mathbf Q]\),
and put \(K_e=K_{2e}(\mathcal O_F)\). At each prime \(\ell\) choose
any finite coefficient extension \(E_\ell/\mathbf Q_\ell\), with
integers \(O_\ell\) and ramification index \(e_\ell\), and any
\(O_\ell^\times\)-valued Dirichlet character \(\psi_\ell\) of
modulus \(N\). Let \(S_\ell\) be the strict kernel with coefficients
\((E_\ell/O_\ell)(\psi_\ell^{-1}\chi_\ell^{-e})\), trivial on every
inertia group and locally trivial at the decomposition groups above
\(\ell N\), in the sense of restriction of cohomology classes.
Each \(S_\ell\) is finite. It vanishes when
\(\ell\nmid n\,\#K_e\). Moreover
\[
 \sum_{\ell}\frac{\log\ell}{e_\ell}\operatorname{length}_{O_\ell}S_\ell
 \ \le\ \log\#K_e+n\log(4Nne).
\]
If one sums over a family of at most \(\varphi(N)\) such characters
at each prime, the right side is multiplied by \(\varphi(N)\).
\end{proposition}
\begin{proof}
Over \(F\) each \(\psi_\ell\) is trivial. On inertia away from \(\ell\)
the remaining cyclotomic twist also acts trivially. The restriction
of a strict cocycle therefore vanishes on each such inertia group,
since a coboundary there is zero. It factors through the quotient
unramified outside \(\ell\), and its class restricts to zero on the
decomposition groups at \(\ell\). There are no real places of \(F\).
Choose a \(\mathbf Z_\ell\)-basis of \(O_\ell\) and put
\(d_\ell=[E_\ell:\mathbf Q_\ell]\). The restricted coefficient module
is a sum of \(d_\ell\) copies of \((\mathbf Q_\ell/\mathbf Z_\ell)(-e)\).
Poitou--Tate identifies the dual of the strict kernel for one copy
with the kernel of localization in
\(H^2(G_{F,S(\ell)},\mathbf Z_\ell(e+1))\), where \(S(\ell)\)
is the set of primes of \(F\) above \(\ell\).
The degree-two injection, localization and Quillen--Lichtenbaum
comparison at the start of the proof of Theorem~\ref{newcoh:lemmaK}
embed this kernel in \(K_e\{\ell\}\). Those steps only require that
\(F\) be totally imaginary and \(e\ge1\). Equivalently, the even
case of Weibel's Theorem~VI.8.4 gives
\(K_e\{\ell\}\simeq H^2_{\mathrm{et}}(\mathcal O_F[1/\ell],
\mathbf Z_\ell(e+1))\), also for \(\ell=2\)~\cite{weibelKbook}.
The root-of-unity assumption entered the subsequent generator count.
Thus the order of the restriction image is at most
\(\#(K_e\{\ell\})^{d_\ell}\).

On the restriction kernel, restriction followed by corestriction
gives multiplication by \(n\), so the kernel vanishes if
\(\ell\nmid n\). For \(\ell\mid n\) write
\(r=\max(1,v_\ell(4N))\); at two this is at least two.
The image of \(G_F\) under \(\chi_\ell\) contains \(1+\ell^r\).
Indeed a cyclotomic automorphism with that \(\ell\)-adic component
and identity at every other prime fixes \(\mu_{4N}\), by the Chinese
remainder theorem. The binomial valuation formula gives
\(v_\ell((1+\ell^r)^e-1)=r+v_\ell(e)\), including at two since
\(r\ge2\) there. The \(G_F\)-invariant subgroup of the original
coefficient module therefore has order at most
\(\ell^{d_\ell(r+v_\ell(e))}\).
By inflation--restriction the remaining kernel injects into the
first cohomology of \(\operatorname{Gal}(F/\mathbf Q)\) with these
invariants. Counting all functions on that group bounds its order
by the \(n\)-th power of the displayed bound. This proves the claimed
finite support since \(K_e\) is finite.

For any finite \(O_\ell\)-module \(M\),
\(\log\#M/d_\ell=(\log\ell/e_\ell)\operatorname{length}_{O_\ell} M\).
After this normalization the restriction image contributes at most
\(\sum_\ell\log\#(K_e\{\ell\})=\log\#K_e\).
The kernel contributes at most \(n\log(4Nne)\), since the product
of the distinct primes dividing \(n\) divides \(n\), the product
of their \(\ell^{v_\ell(4N)}\) divides \(4N\), and their
\(\ell^{v_\ell(e)}\) have product at most \(e\). This estimate
is independent of the character and the coefficient field.
\end{proof}

\begin{corollary}[Summing distinct characters before bounding their lengths]
\label{newcoh:tamepool}
Retain the notation of Proposition~\ref{newcoh:fullsupport}. Regard
the characters at \(\ell\) as distinct characters of
\((\mathbf Z/N\mathbf Z)^\times\) in one algebraic closure of
\(\mathbf Q_\ell\). Put \(b_\ell=\ell^{v_\ell(\varphi(N))}\),
\(a_\ell=\max(1,v_\ell(4N))+v_\ell(e)\), and let \(s_\ell\)
be the least number of generators of the Sylow \(\ell\)-subgroup of
\(G=\operatorname{Gal}(F/\mathbf Q)\), taking zero for the trivial
subgroup. For any subset of these characters one has
\[
 \sum_{\psi}
  \frac{\log\ell}{e_{\ell,\psi}}\,
     \operatorname{length}_{O_{\ell,\psi}} S_{\ell,\psi}
 \le b_\ell\bigl(\log\#K_e\{\ell\}
                     +s_\ell a_\ell\log\ell\bigr).
\]
In particular, at every \(\ell\nmid [F:\mathbf Q]\) the sum over
\emph{all} the characters costs at most \(\log\#K_e\{\ell\}\).
\end{corollary}
\begin{proof}
Enlarge all coefficients to a common finite extension with integers
\(O\), splitting the finite abelian group \(G\). A further extension
of discrete valuation rings is finite free. Continuous cochains and
the kernels of their restriction maps commute with this scalar
extension: after choosing a module basis the underlying operation
is a finite direct sum. The length in the larger integers is the
old length times the relative ramification index, by a composition
series. Thus the normalized lengths are unchanged.

Write \(G=P_\ell\times H\), with \(P_\ell\) its Sylow \(\ell\)-subgroup
and \(H\) of order prime to \(\ell\). Over \(F\) use the single strict
kernel \(M\) with coefficients \((\mathbf Q_\ell/\mathbf Z_\ell)(-e)\)
and localizations at \(\ell\), as in the preceding proof.
It is a finite \(G\)-module of order at most \(\#K_e\{\ell\}\).
Untwisting the restriction of \(S_{\ell,\psi}\) puts its image in
the \(\psi|_H\)-part of \(M\otimes_{\mathbf Z_\ell}O\): the restricted
class was invariant for the action twisted by \(\psi^{-1}\).
These parts for different characters of \(H\) form a direct sum
by the character projectors \( |H|^{-1}\sum_{h\in H}\eta(h)^{-1}h\).
At most \(b_\ell\) of the original characters have a given restriction
to \(H\), since \(G\) surjects onto \((\mathbf Z/N\mathbf Z)^\times\)
and the characters of this quotient trivial on the image of \(H\)
are its \(\ell\)-power order characters. Dividing log orders by
\(d_\ell=[\operatorname{Frac}O:\mathbf Q_\ell]\), the total image
contribution is at most \(b_\ell\log\#K_e\{\ell\}\).

For the kernel put \(U=((\mathbf Q_\ell/\mathbf Z_\ell)(-e))^{G_F}\).
It is cyclic of order at most \(\ell^{a_\ell}\) by the cyclotomic
calculation above. If it is nonzero, \(H\) acts on it through one
prime-to-\(\ell\) character. Indeed the kernel of reduction of the
units of \(\mathbf Z/\ell^j\mathbf Z\) modulo \(\ell\) is an
\(\ell\)-group, and the prime-to-\(\ell\) roots have unique lifts.
Thus \((U\otimes O)(\psi^{-1})^H\), with invariants taken after
twisting, is zero except for at most \(b_\ell\) of the characters.
Inflation--restriction identifies the cohomology in degree one of
\(G\) on this twisted module with that of \(P_\ell\) on its
\(H\)-invariants, since \(|H|\) is invertible. A cocycle on \(P_\ell\)
is determined by its \(s_\ell\) values on generators, by the cocycle
identity. The normalized log order for each of the remaining kernels
is therefore at most \(s_\ell a_\ell\log\ell\). Adding the two
contributions proves the assertion. If \(\ell\nmid n\), then
\(s_\ell=0\) and \(b_\ell=1\) since \(\varphi(N)\mid n\).
\end{proof}

This proves qualitative finiteness for every character under
discussion. To get the required size in the moving windows by this
route one would need a sufficiently sharp uniform bound as \(N,e\)
vary. Finiteness of the single group \(K_e\) for each pair \(N,e\)
alone provides no such rate. Theorem~\ref{famsel:rate} gives an explicit
all-prime bound, and Corollary~\ref{famsel:subwindow} identifies the
part of the required budget it covers.

\subsection{Two vanishing ranges in case III and the integral cotangent inequality}
\label{newcoh:oddbranch}

In case III the conductor of the odd-order character \(\psi\) divides
the fixed target conductor \(q\). Indeed, its component at \(2\) has
both odd order and \(2\)-power order. At either auxiliary prime away
from \(q\) the host condition says that the square of the component
is trivial. At a prime dividing \(q\) the host allows no higher
conductor there than the local factor of \(q\) (the auxiliary primes
occur to exponent one). The Chinese remainder theorem proves the
claim. Thus only finitely many primitive characters occur as the host
varies. The following criterion applies before any conditions at the
auxiliary level are imposed.
\src{src/Development/AutomorphicArithmetic/AutomorphicVarieties/HeightAndVanishing/HostOdd.lean}{70}{odd-order host characters}.

\begin{proposition}\label{newcoh:oddzero}
Let \(\psi\) have conductor \(f\), odd order \(m>1\), and fixed field
\(K=\overline{\mathbf Q}^{\ker\psi}\). Let \(\ell\ge5\) be prime
to \(mf\), let \(e>0\) be even with \(\ell-1\mid2e\), and set
\(s=\operatorname{ord}_m(\ell)\). Let \(O\) be the ring of integers
of a finite extension of \(\mathbf Q_\ell\) containing the values of
\(\psi\), and put \(A=(\operatorname{Frac}O/O)(\psi^{-1}\chi_\ell^{-e})\).
In either of the following situations \(S_\ell(\mathbf Q,A)=0\).
\begin{enumerate}
\item \(\ell-1\mid e\), \(\psi(\ell)\ne1\), and \(\ell^s\nmid h_K\).
\item \(\ell-1\nmid e\) and \(\ell^s\nmid h_{K(\sqrt{\ell})}\).
In particular this holds, for fixed \(K\), for all sufficiently large
\(\ell\) with \(s>m/2\).
\end{enumerate}
Here \(h_L\) denotes the ordinary class number of \(L\).
\end{proposition}
\begin{proof}
Write \(\theta=\psi^{-1}\chi_\ell^{-e}\), let \(\pi\) be a uniformizer,
and set \(k=O/(\pi)\), \(\alpha=\bar\theta\) and
\(L=\overline{\mathbf Q}^{\ker\alpha}\).
Any nonzero strict group contains a nonzero class \(x\) killed by
\(\pi\), by the finite-image argument in Proposition~\ref{newcoh:zero}.
For a representative \(z\) we have \(\pi z=\delta b\). Divide \(b\)
by \(\pi\) in the divisible module \(A\) and subtract the resulting
coboundary from \(z\). This gives a representative valued in
\(A[\pi]\simeq k(\alpha)\).

In the first case \(\alpha\) is nontrivial on \(D_\ell\) by the
Frobenius hypothesis. In the second it is nontrivial on inertia,
since \(\psi\) is unramified there and the residual Tate character
is quadratic. In both cases \(A^{D_\ell}=0\), since some element
acts by a scalar whose difference from one is a unit. On \(D_\ell\)
write \(z=\delta b_\ell\) by strictness. Then \(\pi b_\ell\) is
invariant and \(b_\ell\in A[\pi]\). Consequently \(z\) vanishes on
\(D_\ell\cap G_L\). For \(p\ne\ell\) the scalar on
\(I_p\cap G_L\) is already one in \(O\), since the cyclotomic factor
is unramified and reduction is injective on the \(m\)-th roots of
unity. The inertia condition again gives \(z=0\) there. These
conclusions hold at every prolongation, by the conjugation identity
in Proposition~\ref{newcoh:zero}. The restriction to \(G_L\)
therefore factors through \(\operatorname{Cl}(L)\), since real
inertia has order at most two. It is nonzero by inflation--restriction,
as \([L:\mathbf Q]\mid2m\) is prime to \(\ell\). Its image is an
additive subgroup of \(k\) stable under \(\alpha(G_{\mathbf Q})\)
by the cocycle identity. This cyclic scalar group contains the
\(m\)-th roots of unity, since the two factors have coprime orders.
The image is thus a nonzero vector space over
\(\mathbf F_\ell(\mu_m)\), giving \(\ell^s\mid h_L\).

In the first case \(L=K\). In the second, the residual quadratic
factor cuts out \(\mathbf Q(\sqrt{\ell})\), since even \(e\) forces
\(\ell\equiv1\bmod4\). Coprimality of the character orders gives
\(L=K(\sqrt{\ell})\), of degree \(2m\). The tensor product of
integral bases generates an order of discriminant
\(|D_K|^2\ell^m\) in \(L\). Minkowski supplies an ideal in each
class of norm at most \(y=2^m|D_K|\ell^{m/2}\).
There are at most \(\tau_{2m}(n)\) ideals of norm \(n\), by comparison
of each Euler factor with \((1-p^{-z})^{-2m}\) coefficientwise.
Summing the divisor expansion, with the final variable counted by
its upper bound, gives
\[
 h_L\le y(1+\log y)^{2m-1}.
\]
If \(s>m/2\), this is less than \(\ell^s\) for all sufficiently
large \(\ell\), uniformly in \(e\). For example, cubic characters
give the second range at primes \(\ell\equiv5\bmod12\).
\end{proof}

For the remaining prime and weight combinations the rank-four
argument compares the strict group with integral derivations.
We give the length inequality in the form needed for this comparison.
Let \(O\) be the ring of integers of a finite extension of \(\mathbf Q_\ell\),
\(B=\bigoplus_i B_i\) a graded \(O\)-algebra, and \(x:B\to O\)
an \(O\)-point, \(J=\ker x\), and
\(U=\{b\in B:\exists\,0\ne t\in O,\ tb\in J^2\}\).
Fix \(j\) with \(x(B_j)=0\) and equations \(F_i\in B_j\).
In \(B/U\) let \(C_j\) be the image of \(B_j\), and put
\(Q_j=C_j/\langle\bar F_i\rangle_O\).

\begin{lemma}\label{newcoh:cotangent}
Suppose \(H\) is an \(O\)-module, \(0\ne a\in O\), and its elements
index distinct \(O\)-linear maps \(D_h:B\to O/(a)\) killing all
\(F_i\). Suppose each \(D_h\) lifts to an \(O\)-valued derivation
at \(x\) vanishing on \(B_i\) for every \(i\ne j\). Then
\(\operatorname{length}_O H\le\operatorname{length}_O Q_j\).
\end{lemma}
\begin{proof}
The Leibniz rule kills \(J^2\) under an integral lift, and
torsion-freeness of \(O\) then kills \(U\). Thus each \(D_h\)
induces a functional \(Q_j\to O/(a)\). Its values on \(B_j\)
determine its values on all of \(B\) by the weight assumption,
so these functionals remain distinct. If \(Q_j\) has finite length,
write its elementary divisors as \(O/(\pi^{b_i})\). The dual with
values in \(O/(a)\) has length
\(\sum_i\min(b_i,v_\pi(a))\). Comparing cardinalities now proves
the inequality, since a finite \(O\)-module of length \(b\) has
\(|k|^b\) elements. Infinite right-hand length gives the assertion
automatically.
\end{proof}
\src{src/Development/AutomorphicArithmetic/AutomorphicVarieties/CommutativeAlgebra/GradedModulesLength/Mprime.lean}{38}{the saturated integral cotangent inequality}.

To apply this comparison let \(H\) be the strict group for the
scalar \(\theta\). Fix \(b\ne0\) and work on its subgroup \(H[b]\).
Suppose \(\theta(g_0)-1\) is a unit. Every class in this subgroup
then has a unique representative \(z\) with \(z(g_0)=0\), obtained
by subtracting \(\delta(z(g_0)/(\theta(g_0)-1))\).
Uniqueness gives additivity and \(bz=0\) as a cocycle. On each of
finitely many designated decomposition groups choose \(s_v\) with
\(\theta(s_v)\ne1\). If \(z|_{D_v}=\delta u_v\), the equality
\(bz(s_v)=0\) shows that \(b(\theta(s_v)-1)u_v=0\). Hence for
\[
 a=b\prod_v(\theta(s_v)-1)
\]
all representatives and all local witnesses have values in
\(A[a]\simeq O/(a)\). This yields an additive injection of \(H[b]\)
into finite-coefficient cocycles satisfying the prescribed local
coboundary conditions. The denominator \(a\) may grow with \(b\).
\src{src/Development/AutomorphicArithmetic/AutomorphicVarieties/SymplecticGSp4/ShaCoc.lean}{112}{common denominator for the strict local witnesses}.

Section~\ref{ct:section} constructs the two square-zero representations
and, from the specified relative presentations, sends their cocycles to
reductions of integral weight \(-1\) derivations killing the equations modulo \(a\).
In the lower left matrix entry with indices \((3,1)\) in one-based
notation the derivative recovers \(z(g)\) times a unit for each \(g\).
This makes the assignment injective, so Lemma~\ref{newcoh:cotangent}
bounds the length of \(H[b]\) by that of \(Q_{-1}\), independent
of \(b\). The subsequent geometric inequality is
\(\operatorname{length}_O Q_{-1}\le2\sum_\xi v_\pi(\xi(a_0))\),
where \(\xi\) runs through a finite cover of the relevant orbits
and \(a_0\) interpolates a trace difference. Applying it still
requires the integral deformation model, rigidity of the relevant
irreducible points, and simultaneous nonvanishing of that trace
difference. In particular the extension groups characterized by
Bella\"iche and Chenevier for a fixed Cayley--Hamilton quotient
are the extensions realized by that quotient. Their theorem alone
does not realize every strict class in this rank-four model.
See Bella\"iche and Chenevier, Theorems 1.5.5--1.5.6~\cite{bellaicheChenevier}.
\src{src/Development/AutomorphicArithmetic/AutomorphicVarieties/SymplecticGSp4/MPrime.lean}{102}{weight minus one and the cocycle entry},
\src{src/Development/AutomorphicArithmetic/AutomorphicVarieties/SymplecticGSp4/GaloisRepIneqOfAllIrreducibleNonvanishing.lean}{152}{rigidity and the orbit valuation inequality}.

\subsection{Strict cocycles as integral tangent functionals}
\label{ct:section}

Let \(O\) be the integers in a finite extension of \(\mathbb Q_\ell\),
\(0\ne a\in O\), \(M=O/aO\), and \(\Gamma=G_{\mathbb Q}\).
Write \(O\ltimes M=O\oplus\epsilon M\), with \(\epsilon^2=0\).
This is a complete local \(O\)-algebra. Its quotients
\((O/\pi^h)\ltimes M\), \(h\ge v_\pi(a)\), are Artin rings and their
inverse limit is \(O\ltimes M\). For a continuous augmented
\(O\)-algebra \(B\xrightarrow{x}O\), a continuous map to this
extension over \(x\) has the unique form \(x+\epsilon d\), with
\(d(fg)=x(f)d(g)+x(g)d(f)\).

Use one-based indices for matrix units. The two diagonal points
\(r_0\), their infinitesimal directions \(X\), and the weights \(j\)
of their nonzero entries are as follows:
\[
\begin{array}{c|c|c|c}
 n&\text{matrix of }r_0& X& j\\ \hline
 3&\operatorname{diag}(\Psi,1,\Psi^{-1})&E_{21}-E_{32}&-1\\
 4&\operatorname{diag}(\alpha\Psi,\Psi,\alpha,1)&E_{31}+E_{42}&(0,-1).
\end{array}
\]
In rank three one appends the sign character \(\eta\) to the matrix,
in the semidirect product in which the negative sign acts by
\(g\mapsto Jg^{-t}J\), \(J\) the antidiagonal matrix of ones. Its
differential \(Y\mapsto-JY^tJ\) fixes \(X\).
In rank four the symplectic adjoint sends \(X\) to \(-X\);
the diagonal point has similitude \(\alpha\Psi\).
In both ranks \(\operatorname{Ad}(r_0(g))X=\theta(g)X\) with
\(\theta=\Psi^{-1}\). The rank-four bidegrees of the four basis
vectors are \((1,1),(0,1),(1,0),(0,0)\).

\begin{proposition}[The square-zero representation]\label{ct:point}
Let \(z:\Gamma\to M\) be a continuous cocycle for \(\theta\).
Then \(r_z(g)=(1+\epsilon z(g)X)r_0(g)\) is a continuous
representation over \(O\ltimes M\) with the prescribed sign or
similitude. If on a subgroup \(D\) one has
\(z(g)=(\theta(g)-1)m\), the restriction is the reduction modulo \(a\)
in the second coordinate of an integral dual-number conjugate of
\(r_0|_D\).
\end{proposition}
\begin{proof}
Multiplication adds the two square-zero coefficients as
\(z(g)+\theta(g)z(h)=z(gh)\). In rank four the linear terms in the
similitude cancel because the symplectic adjoint of \(X\) is \(-X\).
Choose a lift \(\widetilde m\in O\) of \(m\). Conjugation by
\(1-\epsilon\widetilde mX\) changes the left infinitesimal coefficient
of \(r_0(g)\) to \(\widetilde m(\theta(g)-1)X\).
The matrix entries are continuous in the adic topology since \(r_0\)
is continuous and \(z\) is locally constant.
\end{proof}

For rank three take the normalized representative of each class in
\(H[b]\), where \(H\) is the strict group and \(0\ne b\in O\). At the
marked element \(g_1\) one has \(\Psi(g_1)^2-1\in O^\times\);
subtracting a unique coboundary makes \(z(g_1)=0\). The normalized
representative has values in \((\operatorname{Frac}(O)/O)[b]\simeq O/bO\).
Here the local group at \(\ell\) contains \(s\) with
\(\Psi(s)^2-1\in O^\times\), so its coboundary witness is also killed
by \(b\). Take \(a=b\). At complex conjugation \(c\), the identities
\(\theta(c)=1,\ c^2=1\) give \(2z(c)=0\), hence \(z(c)=0\) since
\(\ell\ne2\). The cocycle identity then gives \(z(cg_1^{-1}c)=0\).
Consequently the matrices at the three elements entering the two
normalization equations agree with those of the base point.
On the prescribed inertia elements away from \(\ell\) where the
residual representation is trivial, the datum requires \(r_0(g)=1\).
There \(\theta(g)=1\), so local strictness gives \(z(g)=0\).
\src{src/Development/AutomorphicArithmetic/AutomorphicVarieties/UnitaryThreeGroups/ShaInjGen.lean}{145}{normalized rank-three cocycles and their root entries}.

In rank three the local condition at \(\ell\) specifies particular
components of a closed locus. Its first-order verification can be made
directly. Let \(U\) represent unrestricted framed local lifts with the
prescribed sign. Put \(H_L=\ker\eta\cap\ker\tau\), where
\(\Psi|_{I_\ell}=\varepsilon_\ell^e\tau\), \(e\ge1\), with \(\tau\)
of finite order. Let \(I_{\rm ord}\subset U\) be the ideal of the
inertia characteristic-polynomial identities and the ordered products
of three shifted matrices on \(H_L\cap I_\ell\), with successive
characters \(\Psi,1,\Psi^{-1}\). Write \(\Delta=\ell d_1d_2\) for the
product of the two consecutive Frobenius-gap equations used to saturate
this ideal. Let \(I_t\) be the intersection of the minimal primes of
\(\sqrt{I_{\rm ord}:\Delta^\infty}\) contained in the kernel of the
diagonal point. The framed universal property makes this point unique.
\src{src/Development/AutomorphicArithmetic/AutomorphicVarieties/PadicHodgeTheory/IsInertiaInfiniteCyclotomic.lean}{133}{the selected ordinary components}.

\begin{lemma}\label{ct:local}
For \(Y\in M_3(O)\), the local framed representation
\((1+\epsilon Y)r_0(1+\epsilon Y)^{-1}\) over \(O[\epsilon]\),
\(\epsilon^2=0\), kills \(I_t\).
\end{lemma}
\begin{proof}
Let \(d:G_{\mathbb Q_\ell}\to\widehat{\mathbb Z}\to\mathbb Z_\ell\)
be unramified degree, normalized on arithmetic Frobenius. The character
\(u_T(g)=(1+T)^{d(g)}\) is continuous for the \((\pi,T)\)-adic
topology by the binomial expansion. Over the domain \(O[[T,V]]\)
use the diagonal point for \(\Psi u_T\), with the same sign as before,
conjugated by \(1+VY\). The antidiagonal inverse-transpose involution
fixes every matrix \(\operatorname{diag}(b,1,b^{-1})\), so this
does define a representation of the full local group. On
\(H_L\cap I_\ell\) the diagonal characters are unchanged. The ordered
product of the three shifted diagonal matrices is zero since each
diagonal entry has its own vanishing factor. Conjugation preserves this
identity and the three characteristic coefficients.

At the distinguished inertia element \(s_0\in H_L\) the three diagonal
values \(\lambda_i\) are distinct: \(\varepsilon_\ell(s_0)\) has
infinite order. Set \(c_i=\prod_{j\ne i}(\lambda_i-\lambda_j)\).
The product of the other two shifted matrices is \(c_iE_{ii}\) in
the diagonal basis. At the specified \(\phi\in H_L\) of nonzero
unramified degree, put \(b=\Psi(\phi)(1+T)^{d(\phi)}\). Its three
diagonal values \(b_i\) are \(b,1,b^{-1}\). The trace defining the
\(i\)-th projected Frobenius value is \(c_i b_i\), and hence
\[
 d_i=c_i c_{i+1}\bigl(b_i-\varepsilon_\ell(\phi)b_{i+1}\bigr)
 \qquad(i=1,2).
\]
These are invariant under the displayed conjugation since all the
elements in the traces have positive sign. Both are nonzero: \(b\)
has nonzero linear coefficient. Thus the map \(U\to O[[T,V]]\)
kills \(I_{\rm ord}\) and its saturation and radical. Its prime kernel
contains a minimal prime \(P\) of that radical. Evaluation at
\(T=V=0\) is the diagonal point, so \(P\) occurs in the intersection
defining \(I_t\). In particular the map kills \(I_t\). Specializing
\(T=0,\ V=\epsilon\) proves the assertion, including when a gap
vanishes at the original diagonal point.
\end{proof}

In rank four use the common denominator for the local witnesses as
in Lemma~\ref{newcoh:cotangent} and the paragraph following it.
At \(\ell\) the indicated ordinary equations prescribe the
characteristic polynomial at \(s_0\), the vanishing of the forbidden
entries between its polynomial projectors, and the prescribed diagonal
inertia characters between equal projectors. These hold on the
diagonal representation and survive symplectic conjugation, since a
polynomial in a matrix transforms by the same conjugation. At
\(p\in S_b=\{p\in S:p\ne\ell,\ \ell\mid\#\psi(I_p)\}\) the condition
is fixed similitude and triviality on \(I_p\cap\ker\psi\); the diagonal
datum has this property, and conjugation preserves it. At all of these
places Proposition~\ref{ct:point} supplies an integral conjugate using
the local witness. Away from \(T=\{\ell\}\cup S_b\), the datum has
\(r_0(g)=1\) on inertia elements with trivial residual action, and the
strict cocycle vanishes there. Thus the local lifts and global \(r_z\)
agree after reduction, with every comparison frame equal to the identity.
The representing framed global ring consequently gives a map into
\(O\ltimes M\); its restriction to the local coefficient ring lifts
to \(O[\epsilon]/(\epsilon^2)\). The same assertions hold for the
normalized rank-three ring by Lemma~\ref{ct:local} and the checks above.
\src{src/Development/AutomorphicArithmetic/AutomorphicVarieties/SymplecticGSp4/ExistsLocalFirstOrderLifts.lean}{188}{integral local conjugates}
and \src{src/Development/AutomorphicArithmetic/AutomorphicVarieties/SymplecticGSp4/ExistsQuotientDerivationOfShaCocycle.lean}{252}{the global framed point}.

\begin{lemma}[Integral lift on a relative presentation]\label{ct:lift}
Let \(L\) be an augmented complete local \(O\)-algebra with weighted
coordinates at the augmentation, and \(P=L[[t_1,\ldots,t_s]]\) with
homogeneous variables of augmentation zero. Suppose a map
\(P\to O\ltimes M\) over the augmentation restricts on \(L\) to the
reduction of an integral derivation of weight \(j\) in its second
coordinate. Suppose its second coordinate on \(t_i\) is zero when
\(\mathrm{wt}(t_i)\ne j\). Then its second coordinate lifts to an
integral derivation of weight \(j\) on \(P\).
\end{lemma}
\begin{proof}
Lift its value on each \(t_i\) to \(d_i\in O\), taking \(d_i=0\) in
all other weights. If \(D_L\) is the given local lift and \(x_L\) is
the augmentation of \(L\), the required map is
\[
 D\Bigl(\sum_\nu b_\nu t^\nu\Bigr)
     =D_L(b_0)+\sum_i x_L(b_{e_i})d_i .
\]
Comparing constant and linear coefficients proves Leibniz and
reduction. The augmentation is fixed by the grading, so this map
vanishes on the other homogeneous pieces.
\end{proof}

Apply the lemma on the weighted presentation before the global
relations. The presentation variables are shifted matrix entries or
comparison-frame entries. Derivatives of the frames vanish, and the
matrix derivatives are precisely the entries of \(z(g)Xr_0(g)\).
The rank-three presentation also uses the inverse-transpose block:
the derivative of inversion at a diagonal \(D\) is
\(Y\mapsto-D^{-1}YD^{-1}\); transposition and reversal of the indices
preserve the assigned block weight. Thus all these variables satisfy
the hypothesis on the weights. The local matrix-entry derivatives of
the integral conjugates have the same weight. More generally, suppose
maps from the local representing ring into integral dual numbers are
determined by the local matrix entries. For an integral torus unit
\(t\), precompose the local map by the weight action of \(t\), or
multiply its second coordinate by the character \(t^j\). The two maps
agree on all matrix entries and hence everywhere. On a homogeneous
element of weight \(i\) this gives \((t^i-t^j)d(f)=0\). For \(i\ne j\)
one can separate the two weights using \(1+\ell\) in one torus
coordinate. Since \(O\) is a domain this gives \(d(f)=0\). This
argument verifies the weight condition on the whole local derivation.
\src{src/Development/AutomorphicArithmetic/AutomorphicVarieties/UnitaryThreeGroups/ModelOfPres.lean}{302}{the relative rank-three presentation}.

In rank four the presentation is further cut by six homogeneous frame
or matrix entries per distinguished place, each with nonzero first
weight. The integral derivation of weight \((0,-1)\) kills them.
In the resulting quotient \(P'\), the model next divides by the kernel
of its characteristic-zero formal chart. The precise property needed
here is that for \(f\) in this kernel some \(0\ne c\in O\) satisfies
\(cf\in(\ker x_{P'})^2\). Then \(cD(f)=0\), so the integral derivation
descends through the chart kernel. This is the cotangent injectivity
condition of the chart after extension to \(\operatorname{Frac}(O)\).
The other normalization equations involve the frames at \(S_b\) and
off-diagonal entries of \(r(g_0)+\alpha(g_0)\Psi(g_0)r(g_0^{-1})\).
Their derivatives vanish modulo \(a\) because the frames are constant
and \(z(g_0)=z(g_0^{-1})=0\). All original global relations also vanish
modulo \(a\), since \(r_z\) defines a point of the framed ring.
\src{src/Development/AutomorphicArithmetic/AutomorphicVarieties/SymplecticGSp4/ExModel.lean}{496}{passage to the normalized weighted quotient}
and \src{src/Development/AutomorphicArithmetic/AutomorphicVarieties/SymplecticGSp4/ExistsSmoothQuotientByWeightedEquations.lean}{104}{the chart-kernel condition}.

On the sum of homogeneous pieces used in Lemma~\ref{newcoh:cotangent},
restrict the resulting derivation to weight \(-1\) (weight \((0,-1)\)
before taking the first-weight-zero algebra in rank four). It kills
the saturation of the square of the augmentation ideal because it
is integral, and its reduction kills the indicated equations.
The functional on \(Q_{-1}\) recovers \(z(g)\) from a homogeneous
lift of matrix entry \((2,1)\) or \((3,1)\), respectively, by dividing
by \(r_0(g)_{11}\in O^\times\). Distinct normalized cocycles therefore
give distinct functionals. In rank three take \(a=b\). In rank four
take \(a=b\prod_v(\theta(s_v)-1)\), as above, so that each local
witness is represented at this denominator. In either rank the same
\(Q_{-1}\) bounds the length of \(H[b]\) for every nonzero \(b\).
Every continuous cocycle into the discrete torsion module has finite
image by compactness, so \(H=\bigcup_hH[\pi^h]\). If \(Q_{-1}\)
has finite length, this increasing sequence stabilizes since a strict
inclusion would increase length. This proves finiteness and the
length bound for \(H\) in both ranks. The following sections construct the weighted presentations, count their
equations and prove cotangent injectivity at the rank-four cut.
Finiteness of the invariant quotient and rigidity at the remaining
global points are further requirements of Theorem~\ref{newcoh:orbitdegree}.

\subsection{The framed presentation and its real-place defect}
\label{dp:section}

Let \(O\) be the ring of integers of a finite extension \(E/\mathbf Q_\ell\),
with residue field \(k\) of odd characteristic. The two presentations used
here are presentations of \emph{framed} representations. In rank three
the form-preserving group in the two dual blocks has Lie algebra
\(\mathfrak {gl}_3\). In rank four, fixing the multiplier gives
\(\mathfrak {sp}_4\). Write \(M_j\) for the residual Lie algebra of
weight \(j\), with Galois action by the adjoint of the residual point.
The grading records the weight of an entry \(ab\) as \(w_a-w_b\).
In particular the tangent vector that reads that entry is indexed here
by the weight of the entry it reads.

Write \(G=G_{\mathbf Q,S}\) for the Galois group unramified outside the
finite primes \(S\) and infinity. Partition \(S=T\sqcup Q\), with
\(\ell=v_0\in T\), and use the full local decomposition groups \(G_v\)
with their maps to \(G\). At \(v\in T\) a complete local augmented
algebra \(B\) supplies a simultaneous family \(\rho_v\) of local
representations. In rank three \(T=\{\ell\}\) and \(B\) is the
prescribed quotient of the local ring. In rank four
\(T=\{\ell\}\cup S_b\). At \(q\in Q\) representations must kill
\(N_q=I_q\cap\ker\bar r\), where \(I_q\) is inertia. In both
applications \([I_q:N_q]\) is prime to \(\ell\). Let \(r_0\) be the
chosen integral lift. An object over a local \(O\)-algebra \(C\)
consists of \(i:B\to C\), a lift \(r\) with the prescribed form and
multiplier, and congruence-to-identity isometries \(\alpha_v\), with
\[
 \alpha_{v_0}=1,\qquad
 r(s)\alpha_v=\alpha_v\,i(\rho_v(s))\quad(s\in G_v).
\]
All matrix-valued maps are continuous. Representations over finite
Artin \(O\)-algebras use the discrete topology.

\begin{proposition}[Relative generators and relations]\label{dp:relative}
Suppose that \(B\) has a finite weighted power series presentation, the
entries of \(\rho_v\) have the indicated weights, and the base point
has weight zero. The framed ring \(U\) has a presentation over \(B\)
by \(g_j\) variables and at most \(o_j\) relations of weight \(j\).
The following formulas describe both integers. Put \(M=M_j\), write
\(dy(s)=s y-y\), and set
\[
 {\cal T}(M)=\{(z,(y_v)):\ z\in Z^1(G,M),\
 z|_{N_q}=0,\ z|_{G_v}=dy_v,\ y_{v_0}=0\}.
\]
Then \(g_j=\dim_k{\cal T}(M)\). An obstruction cocycle is a continuous
triple \((c,(b_v),(b_q))\), where \(c\in Z^2(G,M)\) and
\(db_w=c|_{G_w}\) for \(w\in S\). Divide by the triples
\[
 \bigl(dh,(h|_{G_v}-dy_v)_v,(h|_{G_q}-e_q)_q\bigr),
 \qquad e_q\in Z^1(G_q,M),\quad e_q|_{N_q}=0.
\]
The dimension of the resulting space \({\cal O}b(M)\) is \(o_j\).
The variables can be chosen among entries of \(r(g)-r_0(g)\) and
\(\alpha_v-1\). Zero equations can be appended in each weight to
give exactly \(o_j\) listed equations.
\end{proposition}

\begin{proof}
For clarity, the usual finite presentation argument starts with finitely
many matrices. The maximal pro-\(\ell\) quotient of \(\ker\bar r\) in
\(G\) is topologically finitely generated. Indeed its Frattini quotient
is dual to \(H^1(\ker\bar r,\mathbf F_\ell)\), which is finite by
finiteness of number-field cohomology with ramification in a finite
set. The kernel of this maximal quotient is characteristic in
\(\ker\bar r\). Every Artin lift factors through the corresponding
quotient of \(G\), since the congruence kernel of a matrix group over
a finite local \(O\)-algebra is an \(\ell\)-group. Choose topological
generators there, adjoining lifts of the finite residual quotient.
Start with \(B\) and matrix-entry variables shifted by their base
values for the generators and frames. Invertibility holds modulo the
maximal ideal and hence in the power series ring. The free profinite
group on the generators maps continuously to its matrix group by
passage through the finite coefficient quotients. Impose its relators,
the form, block, multiplier, inertia and frame equations. These are
closed ideal conditions and give precisely the stated universal
property, first over Artin rings and then over their inverse limits.
The noetherian ring already contains the limits of its ideals.
Every matrix-product summand has the weight of the resulting entry.
The base lift preserves the weight spaces and normalizes the form and
block conditions. Thus these equations respect the matrix weights,
including the relations obtained by profinite limits. This weight
argument works termwise modulo powers of the maximal ideal.

Hold \(B\) fixed at its residual point and evaluate on \(k[\epsilon]\),
\(\epsilon^2=0\). Write \(r=(1+\epsilon z)\bar r\) and
\(\alpha_v=1-\epsilon y_v\). Multiplication is exactly the cocycle
identity and the frame equations are the equations defining
\({\cal T}(M)\). In particular no quotient by global conjugation
occurs. Matrix-entry evaluations span the dual of this finite
dimensional space because they separate all the pairs. Choose a basis
among these evaluations in each weight. With the corresponding
variables \(t_i\), the map \(P=B[[t_1,\ldots,t_g]]\to U\) is
surjective by successively lifting modulo powers of the maximal ideal.
Its kernel \(I\) lies in \(\mathfrak m_B P+\mathfrak m_P^2\), by the
chosen basis. The kernel is closed under taking weight components:
substitution of the homogeneous images is componentwise modulo every
power, and both rings are separated. Therefore \(I/\mathfrak m_PI\)
has a homogeneous basis lifting to homogeneous generators. Finiteness
of this vector space permits only finitely many weights. Nakayama's
lemma gives generation of the full ideal.

For the bound, put \(C=I/\mathfrak m_PI\). In
\(P/\mathfrak m_PI\) the maximal ideal annihilates \(C\). Choose \(h\)
so large that \(C\) meets the \(h\)-th maximal-ideal power trivially.
This follows from Krull intersection and finite dimension of \(C\).
After truncating by that power, the extension onto \(P/(I+\mathfrak
m_P^h)\) therefore still has kernel \(C\). Lift the universal matrices
as functions across this finite extension, preserving their form,
and lift the frames. A form-preserving congruence matrix admits such
lifts: lift its entries and correct the form error to first order
by subtracting half its adjoint-symmetrization error in the kernel.
In the paired blocks one may equivalently lift the first block and
use its inverse transpose in the second. At \(v\in T\) the specified
lift is supplied by the structural map of \(B\). At \(q\in Q\)
there is a local lift killing \(N_q\): \(G_q/N_q\) is an extension
of \(\widehat{\mathbf Z}\) by the finite prime-to-\(\ell\) group
\(I_q/N_q\). Averaging on the finite subgroup and
\(\ell\)-cohomological dimension one of \(\widehat{\mathbf Z}\)
give \(H^2(G_q/N_q,M)=0\), which kills the local multiplication
error. Locally constant lifts of matrices exist since all the rings
in this step are finite.

For a set-theoretic lift \(\widetilde r\), the excess of
\(\widetilde r(s)\widetilde r(t)\widetilde r(st)^{-1}\) over 1
is a Lie-algebra-valued two-cocycle with coefficients in \(C\).
Comparison with the true local lifts, conjugated by the lifted
frames, gives the cochains \(b_w\). Changing the global lifts
adds \(dh\), changing frames adds a local coboundary, and changing
a local lift killing \(N_q\) adds a cocycle vanishing there.
Conversely these changes exhaust the displayed coboundaries. The
frame at \(v_0\) may be restored to 1 by common conjugation. Thus
the class after a linear pushout \(C\to k\) vanishes exactly when
the point lifts over \(B\) to that pushout. The component supported
on \(C_j\) has its matrix values in \(M_j\): choose homogeneous lifts
of the entries through the graded truncation and take the coefficient
of a degree-\(j\) kernel element in each error. For the auxiliary
local corrections one may project cochains onto the matching weight
since the residual action preserves every weight. Equivalently first
push out to a homogeneous kernel line, where only entries of its own
weight admit a nonzero homogeneous correction.

The resulting linear map \(C_j^\vee\to{\cal O}b(M_j)\) is injective.
Indeed a zero obstruction gives \(U\to A\) lifting the truncated
point in the pushed-out extension \(A\). Its composite with \(P\to U\)
and the canonical map \(P\to A\) differ by a \(B\)-derivation into
the kernel \(k\). This derivation kills
\(\mathfrak m_BP+\mathfrak m_P^2\), hence \(I\). On \(I\) the first
map is zero and the second map is exactly the chosen functional.
This proves the bound, including the case of a zero obstruction
space. The padding by zero relations makes no minimality assertion.
\src{src/Development/AutomorphicArithmetic/AutomorphicVarieties/CommutativeAlgebra/LemmaP.lean}{256}{The relative presentation argument}
and \src{src/Development/AutomorphicArithmetic/AutomorphicVarieties/CommutativeAlgebra/FramedTwoCoboundaries.lean}{249}{its precise functor and indexing}.
\end{proof}

\begin{proposition}[The equation count]\label{dp:count}
For a finite \(k[G]\)-module \(M\) on which every \(N_q\) acts
trivially, the above spaces satisfy
\[
 \dim{\cal O}b(M)-\dim{\cal T}(M)=
 \dim M^{c=1}+\dim H^0(G,M^\vee(1)),
\]
where \(c\) is complex conjugation and \(M^\vee=\operatorname{Hom}_k(M,k)\).
\end{proposition}

\begin{proof}
Let \(L_q\subset H^1(G_q,M)\) be the inflated classes from
\(G_q/N_q\) and let \(A^1\) be the product of \(H^1(G_v,M)\)
for \(v\in T\) and \(H^1(G_q,M)/L_q\) for \(q\in Q\).
Write \(r_1:H^1(G,M)\to A^1\) and
\(r_2:H^2(G,M)\to\prod_{w\in S}H^2(G_w,M)\).
Sending a pair in \({\cal T}\) to its global class has image
\(\ker r_1\). For any class there choose the local primitives and
subtract the global coboundary of the primitive at \(v_0\).
Over the zero class the choices are local invariant vectors at
\(T\) modulo the simultaneous global invariant vector, as is also
seen by starting with an arbitrary global primitive. Forgetting the
local cochains in \({\cal O}b\) has image \(\ker r_2\) by definition,
and kernel \(\operatorname{coker}r_1\), by subtracting a global
primitive of \(c\). Thus
\[
 \dim{\cal T}=\dim\ker r_1+\sum_{v\in T}h^0_v-h^0_G,\qquad
 \dim{\cal O}b=\dim\operatorname{coker}r_1+\dim\ker r_2.
\]
For \(q\in Q\), averaging over \(I_q/N_q\) expresses \(L_q\)
as the coinvariants of Frobenius on \(M^{I_q/N_q}\).
The kernel and cokernel of Frobenius minus 1 have the same dimension,
so \(\dim L_q=h_q^0\). The final segment of Poitou--Tate identifies
\(\dim\operatorname{coker}r_2=\dim H^0(G,M^\vee(1))\).
The real place has zero positive-degree cohomology since \(\ell\)
is odd. With \(\chi=h^0-h^1+h^2\), the remaining terms are
\(\chi_G-\sum_{w\in S}\chi_w\). The local Euler characteristics
are \(-\dim M\) at \(\ell\) and zero at all other finite primes.
The global Euler characteristic is \(-\dim M^{c=-1}\).
These formulas give the assertion. For the finite module formulas
see Milne~\cite{milneADT},
I, Theorems 2.8, 4.10 and 5.1. Passing from group orders to
\(k\)-dimensions divides all logarithms by \(\log |k|\).
\end{proof}

In rank three the auxiliary datum supplies \(g\) with
\(\eta(g)=\bar\Psi(g)=1\) and \(\bar\chi_\ell(g)=-1\).
Indeed its field-disjointness condition gives \(h\) in the joint
kernel of the residual cyclotomic and finite characters, outside
\(\ker\eta\). Hence \(\bar\Psi(h)=1\) and \(\eta(h)=-1\).
Since \(\eta(c)=-1\), \(\Psi(c)=1\), take \(g=ch\).
It acts as minus the identity on the twisted dual. On
\(\mathfrak {gl}_3(k)\) the action of \(c\) is negative transpose
across the antidiagonal. Its fixed part in weight zero is the line
\(\operatorname{diag}(a,0,-a)\). In either first off-diagonal the
two entries are negative partners, giving one line, and each outer
corner is its own negative, giving zero. The relative list therefore
has \(g_j+{\bf1}_{|j|\le1}\) equations. The eight normalization
entries have weights \(0,0,1,-1,2,-2,1,-1\). The resulting equation
counts in weights \((-1,-2,1,2,0)\) are
\[
 (g_{-1}+3,\ g_{-2}+1,\ g_1+3,\ g_2+1,\ g_0+3).
\]
\src{src/Development/AutomorphicArithmetic/AutomorphicVarieties/UnitaryThreeGroups/BoxPresOf.lean}{149}{The rank-three framed ring}
and \src{src/Development/AutomorphicArithmetic/AutomorphicVarieties/UnitaryThreeGroups/MopBase.lean}{105}{the normalization entries}.

In rank four put \(u=(1,0)\), \(v=(-1,1)\), \(p=(0,1)\),
\(t=(1,1)\). The positive roots are \(u,v,p,t\). The zero
weight space has dimension two, and each of the eight root spaces has
dimension one. On weight \(j=(a,b)\), \(G\) acts by
\(\bar\alpha^a\bar\Psi^b\) and \(c\) by \((-1)^a\).
For \(j=0,\pm p\) this character differs from \(\bar\chi_\ell\)
at \(c\). For the six roots \(\pm u,\pm v,\pm t\), the specified
\(s_\ell\) has cyclotomic residue 1 and root-character residue
different from 1. Thus the twisted dual has no invariants and
\[
 o_j=g_j+2{\bf1}_{j=0}+{\bf1}_{j=p\ \mathrm{or}\ j=-p}.
\]
There is a separate arithmetic premise for \(g_j=0\) at the six
roots. It is the strict residual class vanishing for the characters
\(\bar\alpha^{\pm1},\bar\beta^{\pm1},\bar\theta^{\pm1}\) imposed by
the good datum. It kills \(\ker r_1\); the local invariants at every
\(v\in T\) vanish because the root character is nontrivial at \(s_v\).
The first dimension formula of the proof now gives \(g_j=0\).
\src{src/Development/AutomorphicArithmetic/AutomorphicVarieties/SymplecticGSp4/NumsZero.lean}{17}{The residual vanishing characters and their use in the count}.

\subsection{The rank-four local chart at the diagonal point}
\label{dp:chart}

Write \(\chi\) for the \(\ell\)-adic cyclotomic character.
The base characters in order are
\[
 (\theta,\Psi,\alpha,1),\qquad
 \Psi=\chi^e\psi,\quad\alpha=\chi\omega^{-1}\xi,\quad
 \beta=\Psi/\alpha,\quad\theta=\alpha\Psi .
\]
The multiplier is \(\theta\). The good datum has \(e\ge2\), finite
characters \(\psi,\xi,\omega\), and \(\xi\) unramified at \(T\).
At \(\ell\) it gives \(s_0\in I_\ell\cap\ker\psi\cap\ker\omega\)
with \(\chi(s_0)\) of infinite order. It also gives
\(\chi(s_\ell)=1\), \(\bar\omega=\bar\chi\), and
\(\bar\xi(s_\ell)\bar\Psi(s_\ell)^i\ne1\) for every integer \(i\).

\begin{proposition}[Local charts of the framed symplectic ring]\label{dp:local}
At the diagonal point, the completed characteristic-zero local ring
for the ordinary projector equations at \(\ell\) is a weighted power
series ring over \(E\). It has two coordinates of each weight
\(0,u,v,p,t\), and one of each weight \(-u,-v,-p,-t\).
For \(q\in S_b\) the corresponding ring for the condition
\(r(I_q\cap\ker\psi)=1\) has two weight-zero coordinates and one
of each of the eight root weights.
\end{proposition}

\begin{proof}
Use finite local Artin \(E\)-algebras \(A\) with residue \(E\),
given their finite-dimensional \(\ell\)-adic vector space topology.
At \(\ell\) write \(d_i=\psi_i(s_0)\) for the four diagonal values.
They are distinct, since their exponents on \(\chi(s_0)\) are
\(e+1,e,1,0\). The source equations specify
\(\det(X-r(s_0))=\prod_i(X-d_i)\). Hence
\[
 P_i=\prod_{k\ne i}\frac{r(s_0)-d_k}{d_i-d_k}
\]
are orthogonal projectors summing to one. Each image is a direct
summand of rank one by reduction to \(E\). The other equations say
\(P_jr(g)P_i=0\) for \(j>i\), and
\(P_i(r(h)-\psi_i(h))P_i=0\) for \(h\in I_\ell\). They specify
stability of the ordered flag and the prescribed inertial action on
its lines. The multiplier pairs lines \(i\) and \(5-i\) perfectly:
pairing any two eigenlines with \(d_i d_j\ne\theta(s_0)\) gives zero,
and nonsingularity of the entire pairing gives the assertion.
Rescaling eigenvectors lifting the original basis now produces a
symplectic basis reducing to that basis. Conversely an upper
triangular lift in such a basis with the given inertia characters
obeys the projector equations, since the partial sums of the
eigenlines of \(r(s_0)\) are its prescribed flag. The flag is
uniquely recovered by the projectors, including over nonreduced \(A\).

In the open cell of the symplectic flag variety at this flag there
is a unique lower unipotent matrix carrying the standard flag to a
given flag. This is Gaussian elimination with diagonal pivots units:
factor a near-identity change of symplectic basis into a lower
unipotent and an upper triangular factor. The involution
\(g\mapsto J^{-1}g^{-t}J\) preserves the two factors, so uniqueness
in \(\mathrm{GL}_4\) makes the factors symplectic. The lower
unipotent Lie algebra has one line in each negative root. Exponential
and logarithm are polynomial inverse maps here in characteristic
zero, since strictly triangular matrices are nilpotent. Thus the
moving-flag parameters give four smooth negative-weight directions.
It remains to consider a lift in the upper triangular group with
standard flag and fixed multiplier.

The diagonal deformations of \(\alpha\) and 1 are independent
unramified twists, determining the other two by the multiplier.
They lift across every small extension \(A'\to A\): use logarithm
on \(1+\mathfrak m_A\), lift the value at Frobenius additively, and
use \(\widehat{\mathbf Z}\to\mathbf Z_\ell\) to define the lifted
continuous additive character. Here and below a small extension has
kernel \(V\) killed by the maximal ideal of \(A'\). With lifted
diagonal fixed, the obstruction to lifting the upper-unipotent part
is in
\[
 H^2(G_{\mathbf Q_\ell},\mathfrak n\otimes_E V),\qquad
 \mathfrak n=E(\alpha)\oplus E(\beta)\oplus E(\Psi)\oplus E(\theta).
\]
Indeed a continuous vector-space section on the nilpotent entries,
in upper-unipotent exponential coordinates, gives a lifted cochain.
The multiplication error takes values in the additive kernel.
Conjugation on that kernel uses only the residual diagonal, so the
four root spaces split as written. A coboundary error can be
removed by subtracting a continuous one-cochain. Local Tate duality
gives \(H^2(G_{\mathbf Q_\ell},E(\gamma))^\vee
=H^0(G_{\mathbf Q_\ell},E(\chi/\gamma))\).
For \(\gamma=\alpha,\beta,\theta\) the residual value at \(s_\ell\)
differs from 1 by the displayed nontriviality with \(i=0,-1,1\),
respectively, while \(\chi(s_\ell)=1\). For \(\gamma=\Psi\) use
\(\chi(s_0)^{e-1}\ne1\). The four obstructions vanish. Every
surjection of finite local Artin \(E\)-algebras factors into small
extensions, giving formal smoothness.

The framed tangent space for a positive root \(\gamma\) consists of
all \(Z^1(G_{\mathbf Q_\ell},E(\gamma))\). By local Euler
characteristic and the preceding \(H^2=0\), its dimension is
\((1+h^0)+(1-h^0)=2\). The two unramified characters give two
zero-weight parameters. Together with the moving flag this gives
fourteen. The \(\ell\)-adic vector-space duality and Euler formula
used here follow from the finite-module theorems by applying them
to the finite quotients of a stable lattice and inverting \(\ell\).
For this passage the inverse systems of finite cohomology groups
are Mittag--Leffler, and every continuous cochain on the compact
powers of the Galois group takes values in a bounded lattice.

At \(q\in S_b\) the effective local group is \(G_q/(I_q\cap\ker\psi)\).
Its inertia quotient \(\Delta\) is finite since \(\ker\psi\) is
open. Its order may be divisible by \(\ell\). Over \(E\), averaging
over \(\Delta\) is exact. Consequently for each weight space \(W\)
of the symplectic adjoint, \(H^2\) of the effective group is zero
and \(H^1\) is the space of Frobenius coinvariants of \(W^\Delta\).
Indeed a stable lattice in \(W^\Delta\) reduces this last calculation
to cohomological dimension one of \(\widehat{\mathbf Z}\) on its
finite \(\ell\)-power quotients. Thus \(h^1=h^0\).
The small-extension multiplication-error argument with smooth
symplectic matrices gives formal smoothness, and its framed tangent
space has dimension \(h^1+\dim W-h^0=\dim W\) weight by weight.

To pass from this calculation to the completion, the integral local
ring \(R\) represents the actual matrix equations: both the
undivided projector expressions and the inertia-kernel conditions
are closed matrix equations in the unrestricted framed ring.
Write \(x:R\to O\) for the diagonal point and \(J=\ker x\).
Then \(\widehat{(R[1/\ell])_{J[1/\ell]}}\) represents precisely the
above Artin \(E\)-valued lifts. For such a lift, its matrix entries
minus their diagonal base values form a bounded subset of the
nilpotent ideal of \(A\) by compactness. Put them into a finite
\(O\)-lattice \(L\) in that ideal. If \(\mathfrak m_A^d=0\), the
order \(O+L+\cdots+L^{d-1}\) is complete local with residue \(k\)
and contains the entire representation. Its finite \(O\)-module
topology is its maximal-ideal topology, since its nilpotent ideal
is killed by the \(d\)-th power. This gives a map from \(R\)
reducing to \(x\), hence a local map on the indicated completion.
Conversely such a map kills \(J^d\), and \(R/J^d\) is a finite
\(O\)-module by the filtration with quotients \(J^i/J^{i+1}\).
It thus induces a continuous lift with bounded entries. Two
choices of order embed in a common one obtained by enlarging \(L\),
so the universal property also gives uniqueness.

Choose homogeneous lifts of a basis of \( (J/J^2)\otimes_O E\)
in \(J[1/\ell]\), using shifted matrix entries. They give a
surjection from a power series ring onto the completion. Formal
smoothness forces it to be an isomorphism: lift the inverse tangent
map successively to the power series ring modulo its maximal-ideal
powers. The resulting composite on the power series ring has
invertible linear term and is an automorphism by recursive
substitution. The coordinates and inverse are homogeneous in each
finite order, by uniqueness of the inverse. This also supplies
cotangent injectivity for the original integral ring after tensoring
with \(E\). The kernel of the map \(R\) to this completion is the
kernel of \(R\to(R[1/\ell])_{J[1/\ell]}\) by Krull intersection.
It is a finitely generated ideal, so a single \(y\in R\) with
\(x(y)\ne0\) annihilates it, by clearing the finitely many
denominators separately and taking their product.
\src{src/Development/AutomorphicArithmetic/AutomorphicVarieties/PadicHodgeTheory/SpFlag.lean}{309}{The inertia projector equations},
\src{src/Development/AutomorphicArithmetic/AutomorphicVarieties/PadicHodgeTheory/SpLiftL.lean}{32}{the four character checks}
and \src{src/Development/AutomorphicArithmetic/AutomorphicVarieties/SymplecticGSp4/LocBase.lean}{142}{the integral chart requirements}.
\end{proof}

\subsection{Cotangent faithfulness after the symplectic cuts}
\label{dp:cut}

\begin{lemma}\label{dp:cutlemma}
Let \(P\) be an augmented noetherian complete local \(O\)-algebra,
with augmentation ideal \(J\). Suppose a weighted chart
\(\phi:P\to E[[Y_1,\ldots,Y_n]]\) induces an isomorphism from
\((J/J^2)\otimes_O E\) to the cotangent at \(Y=0\).
Let \(u_1,\ldots,u_h\in J\) be homogeneous with independent images in
this rational cotangent. There is a chart of \(P'=P/(u_i)\) with
one fewer variable of the given weight for each cut, again inducing
a cotangent isomorphism. If \(P'\) is further divided by the kernel
of this latter chart, every \(O\)-valued point derivation at its
augmentation descends to the resulting domain.
\end{lemma}

\begin{proof}
Complete the linearly independent \(d\phi(u_i)\) to a basis in each
weight. Replace the first coordinates by the full series
\(\phi(u_i)\), keeping the complementary homogeneous linear
coordinates. The substitution has zero constant terms and invertible
linear part. Its inverse is obtained recursively in total degree.
Each finite degree respects the weights even when negative and zero
weights occur. Set the first \(h\) new coordinates to zero. On
cotangents this kills precisely the span of \(du_i\), and the
cotangent of \(P'\) is \(J/(J^2+(u_i))\). Exact scalar extension
therefore gives the asserted isomorphism. For \(a\) in the new
chart kernel its augmentation vanishes and its differential is zero
after inverting \(\ell\). Some \(0\ne c\in O\) then satisfies
\(ca\in(\ker x_{P'})^2\). Every integral point derivation kills this
square, hence kills \(a\) by torsion-freeness of \(O\). The chart
target is a domain.
\end{proof}

For the application take \(P=B[[t_1,\ldots,t_g]]\) from
Proposition~\ref{dp:relative}. The simultaneous local ring \(B\)
represents the product of the finitely many local functors. Its
rational completion is the power series ring on the union of the
local coordinates by Proposition~\ref{dp:local}. The relative
variables add their own free summands to the cotangent, so \(P\)
has the required chart. The global relations cut out \(U\) inside
this ambient formal spectrum. The next cuts are lifted to \(P\)
from the following entries of \(U\), in one-based indices:
\[
 {\cal P}=\big((1,2),(1,4),(3,2),(2,1),(4,1),(2,3)\big).
\]
At \(\ell\) take these entries of \(r(\phi_0)\), where
\(\phi_0=\lim_n s_\ell^{\ell^{n!}}\) is the element specified by the
datum. At every other \(v\in T\) take these entries of \(\alpha_v\).
All have zero augmentation. Use the six Lie algebra elements,
in the same order,
\[
 E_{12}-E_{34},\quad E_{14},\quad E_{32},\quad
 E_{21}-E_{43},\quad E_{41},\quad E_{23}.
\]
Each lies in \(\mathfrak {sp}_4(O)\) for the antidiagonal pairing with
signs \(1,1,-1,-1\), and has just one nonzero entry among \({\cal P}\).
For a selected element \(X\) conjugate \(r_0\) and all the local lifts
by \(1+\epsilon X\) over \(O[\epsilon]/(\epsilon^2)\), leaving all
frames equal to 1. This preserves the local equations by conjugation.
At \(\ell\) the only nonzero selected derivative is
\(\psi_j(\phi_0)-\psi_i(\phi_0)\). Indeed the residual diagonal at
\(\phi_0\) is the same as at \(s_\ell\), since the values belong to
\(k^\times\) and the residue degree \([k:\mathbf F_\ell]\) divides
\(n!\) eventually. Each pair has different first weight, so
\(\bar\xi(s_\ell)\bar\Psi(s_\ell)^m\ne1\) for all \(m\) makes
the residual entries distinct. At \(v\ne\ell\) keep the global
lift fixed, move only \(\alpha_v\) to \(1+\epsilon X\), and
conjugate the corresponding local diagonal lift by its inverse.
This tests one entry with derivative 1. These point derivations of
\(U\) pull back to \(P\), proving independence for any homogeneous
lifts of all the indicated cuts there.

Lemma~\ref{dp:cutlemma} now gives the post-cut chart and the
integral descent of derivations needed in
Section~\ref{dp:section}'s framed construction. Let \(b=|S_b|\).
In each odd root the cut removes \(b+1\) variables. The remaining
dimensions there are 1 for \(u,v,t\) and 0 for their negatives,
using the residual vanishings above. In weights \((-p,p,0)\)
the dimensions are
\[
 (1+b+g_{-p},\quad 2+b+g_p,\quad 2+2b+g_0).
\]
The remaining frame equations at each bad prime have weights
\(p,-p,0,0\), and the two normalization entries of
\(r(g_0)+\theta(g_0)r(g_0^{-1})\) have weights \(p,-p\).
Adding their counts to Proposition~\ref{dp:count} gives final lists
of lengths \(2+b+g_{-p},\,2+b+g_p,\,2+2b+g_0\), respectively.
They define the quotient of the resulting domain used by the model.
\src{src/Development/AutomorphicArithmetic/AutomorphicVarieties/SymplecticGSp4/ExModel.lean}{155}{The entries and six root vectors}
and \src{src/Development/AutomorphicArithmetic/AutomorphicVarieties/SymplecticGSp4/ExModel.lean}{425}{the normalized model}.

\subsection{The selected rank-three branch at a resonant character}
\label{dp:ordinary3}

Here \(G=G_{\mathbf Q_\ell}\), \(I\) is its inertia, and \(\eta\)
is an unramified character of order dividing two. Put \(H=\ker\eta\)
and \(I_0=I\cap H\cap\ker\tau\), where \(\tau\) has finite image.
The source specifies \(\Psi|_I=\chi^e\tau|_I\) with \(e\ge1\).
The matrix group is \(\mathrm{GL}_3\rtimes\{1,-1\}\), where
\(-1\) sends \(A\) to \(J A^{-t}J\) and \(J\) reverses the basis.
Put \((d_1,d_2,d_3)=(\Psi,1,\Psi^{-1})\). The base point of
the unrestricted local framed ring \(R\) is
\(r_0(g)=(\operatorname{diag}(d_i(g)),\eta(g))\), and its
augmentation is \(x\). For \(s\in I_0\) put \(A_i(s)=r(s)-d_i(s)\)
in the matrix block. The ordinary ideal \(I_{\rm ord}\) fixes the
characteristic polynomial on \(I_0\) to \(\prod_i(X-d_i(s))\)
and kills all entries of \(A_1(s)A_2(s')A_3(s'')\) for independent
\(s,s',s''\in I_0\). The definition then saturates by a specific
element \(\delta\), described below. Write
\[
 I_Y=(I_{\rm ord}:\delta^\infty),\qquad
 I_t=\bigcap_{\substack{\mathfrak p\text{ minimal over }I_Y\\
                          \mathfrak p\subseteq\ker x}}\mathfrak p .
\]
Taking the radical of \(I_Y\) before selecting these primes gives
the same intersection. The framed universal property makes \(x\)
the unique map with this specified matrix point.
\src{src/Development/AutomorphicArithmetic/AutomorphicVarieties/PadicHodgeTheory/IsInertiaInfiniteCyclotomic.lean}{97}{The ordinary equations, saturation and selected primes}.

\begin{proposition}\label{dp:rank3chart}
There is a unique selected prime. The completion of
\((R/I_t)[1/\ell]\) at \(x\) is a power series ring over \(E\)
with coordinate counts \(1,2,3,4,2\) in weights \(-2,-1,0,1,2\).
The kernel of the resulting chart from \(R\) is \(I_t\).
These statements hold also when a ratio of two adjacent diagonal
characters is cyclotomic.
\end{proposition}

\begin{proof}
As in Proposition~\ref{dp:local}, use continuous lifts over Artin
\(E\)-algebras, now of \(r_0\). Choose the specified
\(s_0\in I_0\) with \(\chi(s_0)\) of infinite order.
The eigenvalues \(d_i(s_0)\) are distinct. The characteristic
polynomial gives three rank-one projectors. Write \(V_1\) for the
first eigenline and \(V_2\) for the sum of the first two.
In eigen-coordinates, \(A_2(s_0)A_3(s_0)\) is a unit multiple of
the first projector, and \(A_1(s_0)A_2(s_0)\) is a unit multiple
of the third. The triple products therefore say that \(r(s)\) acts by
\(d_1(s)\) on \(V_1\), preserves \(V_2\), and acts on the last
quotient by \(d_3(s)\). Subtracting these two diagonal terms from
the specified trace gives \(d_2(s)\) on \(V_2/V_1\). Conversely
these three conditions imply every triple product, by applying its
factors from right to left along this flag.

The flag is unique by its ordered characters at \(s_0\). Conjugation
by \(g\in G\) preserves \(I_0\) and all the \(d_i\). When
\(\eta(g)=-1\) the induced action on flags first takes the dual flag
in reversed order via \(J\); its ordered quotient characters on
inertia are again \(d_1,d_2,d_3\) since \(d_{4-i}^{-1}=d_i\).
Uniqueness therefore gives stability under the full semidirect
action. By lower-unipotent change of basis the moving flag contributes
three free coordinates, two of weight \(-1\) and one of weight
\(-2\). Once the flag is standard all matrix blocks are upper
triangular, since inverse transpose across the antidiagonal preserves
the standard flag. Conjugation here and below is in the semidirect
group, so a change of basis also changes the block of an element
outside \(H\) by the prescribed involution.

The diagonal part relative to \(r_0\) is an unramified cocycle in
the nilpotent diagonal torus: on inertia the action on the torus
is trivial, and a character trivial on \(I_0\) factors there
through a finite group with no nonzero additive image in an
\(E\)-algebra. Take logarithms. The three coordinates \(b_i\) of
the value at a chosen Frobenius are free, even if \(\eta\ne1\).
Indeed a cocycle on \(\widehat{\mathbf Z}\) with Frobenius action
\(D=1\) or \(D(b_1,b_2,b_3)=(-b_3,-b_2,-b_1)\) has value
\(\sum_{k=0}^{n-1}D^kb\) at its \(n\)-th power. This extends
continuously using parity and the projection to \(\mathbf Z_\ell\).
Thus the base for the upper unipotent part is
\(\Lambda=E[[b_1,b_2,b_3]]\).

First divide the upper unipotent group by its central line \(E_{13}\).
The two adjacent entries form an additive module. For fixed diagonal
they are precisely continuous one-cocycles, with no quotient by
coboundaries. If \(H=G\), they give two independent rank-one
problems with respective parameters \(t=b_1-b_2,b_2-b_3\).
Their characters are \(\Psi\exp(t\nu)\), where the continuous
unramified additive character \(\nu\) is 1 on Frobenius.
If \([G:H]=2\), the two adjacent lines are exchanged by the
antidiagonal inverse transpose on their Lie algebra. The module
is induced from the first line on \(H\): this line and its translate
by a Frobenius form a basis, also for the deformed diagonal over
every Artin algebra. Its character on \(H\) is
\(\Psi|_H\exp(t\nu_H)\), with \(t=b_1-b_3\) and \(\nu_H\) equal
to 1 on Frobenius of \(H\). Indeed the value of the unramified
diagonal logarithm on the square of Frobenius is \(b+Db\).
This description by induction works functorially in the diagonal
parameters. In particular at \(t=0\) the indicated module is
isomorphic over the whole diagonal base to the constant module.

Here is the local equation for either one of these problems. Let
\(G_*=G\) in the split case and \(G_*=H\) in the induced case.
Denote the corresponding \(G\)-module at the origin by \(W\) and
put \(m=\dim_E Z^1(G,W)\). For its cocycle functor over \(\Lambda\),
a minimal presentation uses \(m\) variables and at most
\(\dim_E H^2(G,W)\) relations. To justify this version over \(E\),
the functor preserves fibre products of Artin algebras since its
objects are functions subject to equalities, and its tangent is
finite-dimensional by local Euler characteristic. Schlessinger's
prorepresentability criterion~\cite[Theorem 2.11]{schlessingerArtin}
therefore applies (the functor at the residue field is
one point). One can include the three diagonal
parameters in the functor and then work relative to their power
series ring. The small-extension argument in the proof of
Proposition~\ref{dp:relative} applies verbatim using continuous
linear sections of surjections of finite-dimensional \(E\)-spaces;
here the single obstruction is the coboundary of a lifted one-cochain.
The relation bound concerns the actual ideal modulo its maximal-ideal
multiple. Local Tate duality, with Shapiro's lemma in the induced
case, gives
\[
 \dim_E H^2(G,W)=
 {\bf1}_{\,\Psi|_{G_*}=\chi|_{G_*}} .
\]
If this is zero the cocycle algebra is a power series ring.

If the dimension is one, its possible single equation is, up to a
change of variables, exactly \(tY=0\). We verify its order of
vanishing. Choose the relative variables \(X_1,\ldots,X_m\) among
linear evaluations of the cocycle; they have weight one for scalar
multiplication of the cocycle. Modulo \(t\) the module is constant
up to the functorial basis isomorphism just proved, so every cocycle
is an arbitrary linear combination of a basis of \(Z^1(G,W)\) with
coefficients in the maximal ideal of the test algebra. Its completed
algebra over \(\Lambda/(t)\) is free in the \(m\) variables. The
chosen evaluations still give coordinates there by the inverse
function theorem. It follows that every equation is divisible by
\(t\). No relation has an \(X\)-independent term, since zero is
a cocycle for every diagonal.

Vary \(t\) to \(h\) and the cocycle to \(kz\) over
\(E[h,k]/(h^2,k^2,hk)\), choosing any linear diagonal section of
the parameter \(t\). The obstruction to lifting to
\(E[h,k]/(h^2,k^2)\), in the \(hk\) coefficient, maps to
\(\nu_{G_*}\smile [z_*]\) under the preceding cohomology
identification. In the split case this follows by expanding the
cocycle identity with character \(\chi(1+h\nu)\): its error on
\((g,g')\) is \(hk\,\nu(g)\chi(g)z(g')\) up to the sign of the
coboundary convention. In the induced case the cochain map to
\(G_*\) restricts and projects to the inducing line; it commutes
with reduction for the induced basis and gives the identical
calculation. Every class of \(H^1(G_*,E(\chi))\) occurs by
Shapiro. The local Tate pairing of this space with \(H^1(G_*,E)\)
is perfect and \(\nu_{G_*}\ne0\). Hence for some \(z\) the
obstruction is nonzero.

In particular there is one relation, with nonzero mixed term
\(t\sum_i c_i X_i\), some \(c_i\ne0\). It may be chosen
homogeneous of weight one. Indeed the ideal is stable under scalar
multiplication of all \(X_i\), which fixes \(b\). Its weight
components belong to the ideal by polynomial interpolation in each
finite truncation over the infinite field \(E\). All relations lie
in \(t(X_1,\ldots,X_m)\), so a relation with this nonzero
quadratic term generates the one-dimensional space of relations
modulo the maximal-ideal multiple. Its weight-one component still
generates. Write it \(t\sum_i a_i(b)X_i\); some \(a_i\) is a
unit. Use \(Y=\sum_i a_i(b)X_i\) as one of the relative
coordinates. This proves the claimed equation including its
multiplicity.

Finally restore the central root. Its Galois action on \(I_0\)
is \(\Psi^2=\chi^{2e}\), hence its \(H^2\) vanishes by local
duality since \(2e\ne1\). Multiplication errors for a lifted
central entry in a small extension are two-cocycles for this root
module, as above. The map to the first-layer functor is therefore
formally smooth with relative tangent its space of one-cocycles.
This adds free variables also over the singular first-layer ring:
choose relative cotangent coordinates and split the corresponding
power series surjection order by order by formal smoothness.
Equivalently the small-extension relation bound with zero obstruction
gives zero relative relations. The central coordinate has weight two.

We compute the effect of the prescribed saturation in these
coordinates. The source chooses \(f\in H\cap\ker\tau\) mapping to
a positive power of Frobenius. For \(i=1,2,3\) put
\(c_i=\prod_{k\ne i}(d_i(s_0)-d_k(s_0))\in E^\times\).
Its trace expression is \(c_i\,\mathrm{tr}(P_i r(f))\), where
\(P_i\) is the corresponding eigenprojector of \(r(s_0)\). In an
upper triangular basis this equals \(c_i\kappa_i(f)\), with
\(\kappa_i\) the \(i\)-th diagonal character on \(H\). Indeed each
\(P_i\) is upper triangular with only its \(i\)-th diagonal entry
equal to 1. Thus \(\delta\) is, up to a unit in the characteristic-zero
completion, precisely
\[
  \prod_{i=1}^2\bigl(\kappa_i(f)/\kappa_{i+1}(f)-\chi(f)\bigr).
\]
For a resonant adjacent block the corresponding factor is
\(\chi(f)(\exp(a t)-1)\), with \(a=\nu_{G_*}(f)>0\).
It is a unit times \(t\). This formula also covers the two equal
ratios on \(H\) in the induced case. In the quotient by all the
coordinates \(Y\) belonging to resonant blocks, the ring is a
power series domain and the image of \(\delta\) is nonzero:
each adjacent ratio still varies as
\(\Psi(f)\exp(a t)\) for a nonzero diagonal parameter \(t\).
Since \(tY=0\) and \(t\) divides \(\delta\) in each resonant
block, the \(\delta\)-power torsion is exactly the ideal of these
\(Y\)'s. This establishes the calculation even if \(\delta(x)=0\).

For any \(G\)-module \(W\) in an adjacent block the local Euler
formula over \(\mathbf Q_\ell\) gives
\(\dim Z^1(G,W)=2\dim W+\dim H^2(G,W)\).
Saturation removes exactly one variable for each of the latter
dimensions. The adjacent roots thus leave four weight-one
coordinates in all; the central root adds two of weight two.
The three unramified diagonal coordinates and the three lower flag
coordinates give the other asserted counts.

Localization at the characteristic-zero point and completion are
flat and commute with \((I_{\rm ord}:\delta^n)\): take the kernel
of multiplication by \(\delta^n\) on the quotient. The chain
stabilizes by noetherianity. It follows that the completion of the
localization modulo \(I_Y\) is the power series domain just computed.
Faithful flatness makes the localized quotient a domain.
Consequently there is exactly one minimal prime over \(I_Y\)
contained in \(\ker x\), the contraction of the zero ideal of this
local quotient. Its contraction to \(R\) is also the kernel of the
completed chart. This proves the assertion about the selected
integral ideal \(I_t\).
\src{src/Development/AutomorphicArithmetic/AutomorphicVarieties/PadicHodgeTheory/RtChart.lean}{41}{The corresponding twelve-coordinate conclusion}.
\end{proof}

The preceding computation is a statement about the ordinary matrix
ideal in the completed local ring. Application to the integral
presentation uses a weight grading on \(R/I_t\). The unrestricted
matrix ring has homogeneous coordinates given by shifted entries.
The ordinary equations are weight-homogeneous, \(\delta\) has
weight zero, and the map to the completed chart preserves weights.
Its kernel \(I_t\) is thus closed under projection onto each weight.
Conjugation by a matrix congruent to 1 also preserves the selected
branch over arbitrary integral test algebras. Indeed conjugate the
universal \(R/I_t\)-valued family by the generic matrix \(1+X\)
in \((R/I_t)[[X_{ij}]]\). Embed this ring coefficientwise into
\(E[[Y_1,\ldots,Y_{12},X_{ij}]]\) by the chart. The ordinary
equations are conjugation invariant; \(\delta\) is unchanged because
its defining traces use elements of \(H\), and it is nonzero in this
domain. Hence the conjugate annihilates \(I_Y\). For \(a\in I_t\)
some \(s\) with \(x(s)\ne0\) satisfies \(sa\in I_Y\), by the
calculation of the localized prime. The conjugate image of \(s\) has
constant coefficient \(x(s)\), so \(a\) also maps to zero.
Injectivity then gives the same conclusion before passage to \(E\).
Substituting any matrix of elements of the maximal ideal for \(X\)
proves the integral assertion. This supplies the local presentation required by
Proposition~\ref{dp:relative} in rank three. Identifying potentially
crystalline points of the required type on the entire model is
a further arithmetic statement.

In rank four the local and relative arguments use the good-datum
hypotheses, including the separate residual vanishings.
In both ranks, finiteness of the integral invariant quotient,
exclusion of extraneous one-sided points and rigidity of the
nontrivial two-sided points still require arithmetic control of
representations away from the diagonal base point. The local charts
and equation counts alone yield none of these global-point assertions.

\subsection{The special fibre of the invariant algebra}\label{dp:finitecurves}

There is a separate arithmetic test for the finiteness hypothesis in
Theorem~\ref{newcoh:orbitdegree}. In rank three it leaves only
strongly irreducible three-dimensional specializations on \(G_K\).
In rank four it leaves an absolutely irreducible four-dimensional
specialization or an even absolutely irreducible two-dimensional
constituent. The fields of coefficients in these statements have
characteristic \(\ell\).

\begin{lemma}[Detecting a nonfinite invariant algebra]\label{dp:detect}
Let \(O\) be a complete DVR with finite residue field \(k\) and
uniformizer \(\pi\). Let \(S\) be a quotient of a finite formal
power series algebra over \(O\) by a closed ideal stable under
projection onto the weight components for integral vector weights.
Write \(C\subset S\) for the image of the weight-zero power series.
Assume the variables are in the maximal ideal. Then \(C\) is finite
over \(O\) if and only if every continuous local map
\[
               S\longrightarrow k'[[t]],\qquad \pi\longmapsto 0,
\]
for finite extensions \(k'/k\), takes \(C\) into the constants.
\end{lemma}
\begin{proof}
The nonzero minimal exponent vectors of weight zero are finite
by Dickson's lemma, and they generate the monoid of all balanced
exponents by subtraction. Denote the corresponding monomials by
\(m_j\). Substituting these into a formal power series over \(O\)
surjects onto \(C\): fix one factorization for each exponent and
group the terms accordingly. The bounded total degrees of the
\(m_j\) ensure convergence in both descriptions. Thus \(C\) is
complete noetherian local, with maximal ideal \((\pi,m_j)\).
Weight-zero projection gives \(C\cap\pi S=\pi C\).
If \(C/\pi C\) is artinian, a power of its maximal ideal is zero.
Consequently the \(\pi\)-adic and maximal-ideal topologies on \(C\)
are equivalent; lifting a \(k\)-basis of \(C/\pi C\) and expanding
successively in powers of \(\pi\) gives finitely many \(O\)-module
generators. Conversely a finite algebra has artinian special fibre,
whose image in the domain \(k'[[t]]\) is a finite field.

Otherwise some \(m_j\) is not nilpotent modulo \(\pi C\). By the
projection identity it survives on a prime quotient of \(S/\pi S\).
Cut this complete local domain to dimension one while preserving
the image \(b\ne0\) of \(m_j\). At each step of dimension \(>1\)
choose an element of the maximal ideal outside the finitely many
height-one primes dividing \(b\), and take a minimal prime over
that nonzero element. This prime has height one, avoids \(b\), and
is strictly below the maximal ideal. The resulting one-dimensional
complete domain has finite normalization since it is excellent.
This normalization is local by completeness and the domain property,
and is an equicharacteristic complete DVR with finite residue
extension. Its coefficient field is the unique lift of this finite
field, so a uniformizer identifies it with \(k'[[t]]\) over \(k\).
Here \(b\) has strictly positive valuation. This proves the criterion.
\end{proof}

In the two presentations one may use homogeneous coordinates taken
from matrix entries minus their residual diagonal lifts. Indeed these
entries generate the representing ring topologically, by the framed
universal property, and homogeneous cotangent generators suffice by
successive approximation. Equivalently, on the post-cut rank-four
model the entries generate \(S\) in the sense that they distinguish
its maps to complete local algebras. They then span its relative
cotangent space by testing the dual numbers of the residue field.
For these coordinates \(C\) is the closure of the algebra generated
by weight-zero products of entries. Closure does not change module
finiteness: a finite \(O\)-submodule is compact, hence closed, and
an \(O\)-submodule of a finite \(O\)-module is finite. Thus
Lemma~\ref{dp:detect} applies to the actual weight-zero entry algebra.
The image of the invariant part \(B\) of the presentation in the
post-cut ring is \(B/(I\cap B)\). In rank four one first takes
weight zero for the first coordinate and subsequently for the second;
this is simultaneous weight \((0,0)\).
\src{src/Development/AutomorphicArithmetic/AutomorphicVarieties/CurveDVRSpecialization/FinCrit.lean}{122}{Curve detection in rank three};
\src{src/Development/AutomorphicArithmetic/AutomorphicVarieties/SymplecticGSp4/FgMapMk.lean}{24}{From the rank-four weight-zero entry algebra to the invariant quotient}.

Here are two arithmetic facts about such curves. We may enlarge
\(k'((t))\) by a finite extension and normalize its valuation ring
to define constituent spaces, their flags and stable lattices.
Take intersections and saturated quotients of the original lattice
to get stable lattices on subquotients. Their residual constituents
together give the residual semisimplification of the original
representation, by reduction of these exact lattice sequences.

First, a continuous character of a number-field Galois group into
the units of an equal-characteristic DVR with finite residue field
is constant if its inertia characters at every finite place have
finite image, with ramification at only finitely many places.
Indeed a prime-to-\(\ell\) root of unity is the unique lift of its
residue and the principal units have no torsion in characteristic
\(\ell\). Dividing by the lift of the residual character leaves an
everywhere unramified character into principal units. It factors
through the finite narrow class group by global class field theory,
and is therefore trivial. This argument applies on every finite
number-field extension as well. If a matrix representation becomes
triangular there, and its inertial eigencharacters have finite
image, all its diagonal characters there are constant. On the full
group every eigenvalue is then a root of unity, by raising a group
element to a positive power in the subgroup. Its characteristic
polynomial is the constant lift of its reduction, even if the
subgroup and power depend on the curve.

\begin{lemma}[Odd two-dimensional subquotients on a curve]\label{dp:pan}
Let \(\ell>2\). Suppose a continuous two-dimensional representation
of \(G_{\mathbb Q}\) over a finite extension of \(k'((t))\) has a
stable lattice, is unramified outside finitely many primes, has
finite inertial eigencharacters everywhere, and is triangular on
\(G_{\mathbb Q_\ell}\) over an algebraic closure. Suppose its
residual semisimplification is a sum of two characters of opposite
parity. Then its characteristic polynomials agree with the constant
lifts of their reductions; in particular it is geometrically
reducible.
\end{lemma}
\begin{proof}
After a finite coefficient extension twist by the constant inverse
of the first residual character. Write \(1,\bar\chi\) for the
resulting residual characters. By the preceding character argument
the determinant is the constant lift of \(\bar\chi\). Use its
Teichm\"uller lift \(\chi\) to \(W(k')\), enlarging \(k'\) to the
residue field of the chosen extension. It is an odd finite-order
character and is de Rham locally at \(\ell\). We use the number
field \(\mathbf Q\), an abelian totally real field unramified at
\(\ell\), so the required extension of the residual determinant
to \(G_{\mathbf Q}\) is already given. Choose a finite set \(S\)
of finite primes containing \(\ell\) and every ramified prime
of the twisted representation. Pan's ordinary
pseudodeformation ring for \(G_{\mathbb Q,S}\) at \(1+\bar\chi\),
with determinant \(\chi\), is finite over
\(\Lambda=W(k')[[\mathbb Z_\ell^\times(\ell)]]\), where the group
in brackets is the pro-\(\ell\) part of the local units.
Use the distinguished theorem when
\(\bar\chi|_{G_{\mathbb Q_\ell}}\ne1\). If its restriction is
trivial, use the character-labelled theorem and the actual
subcharacter \(\lambda_1\) of the local triangular representation.
Both local diagonal characters are integral units, since the
inverses are integral too. In a triangular basis
\((r(\sigma)-\lambda_1(\sigma))(r(\tau)-\lambda_2(\tau))=0\)
for local \(\sigma,\tau\); taking the trace after multiplying by
any \(r(g)\) gives the extra ordinary pseudocharacter relation.
The coefficient homomorphism from \(W(k')\) to the normalized
curve DVR is reduction to \(k'\) followed by its constant inclusion.
Thus the trace and, in the second case, the chosen character
define the asserted map of universal rings on passing to the
inverse limit of the Artin quotients. The finiteness assertions
used here hold for every odd prime, including three.
The map on \(\Lambda\) comes from an inertia
character which has finite image in the residue-field units, so
it kills the augmentation of the pro-\(\ell\) group and also
\(\ell\). The image of the finite algebra in this DVR is a finite
field. The trace and determinant therefore equal their residual
values as constants. The characteristic-polynomial criterion for
semisimplification (Brauer--Nesbitt) finishes the proof.
\cite[Theorems~5.1.1(1) and~6.1.1(1) in the published
version, Theorems~5.1.2(1) and~6.1.2(1) on pp.~71 and~95--96
of the author manuscript]{panFM}
\end{proof}

We explain how traces control the marked entries required here. Put
\(\mathcal E=\operatorname{span}\{r(g)\}\) over the algebraic
closure of the fraction field of the curve, using \(G_K\) in rank
three. It is an algebra. It contains the individual coordinate
projectors \(E_{ii}\) in both models. In rank three the marked
difference \(Z=r(g_1)-Jr(g_1)^tJ\) belongs to this algebra since
the second matrix is \(r(cg_1^{-1}c)\). The normalization kills
entries \(01,10\), and antisymmetry under reverse transpose makes
\(Z=\operatorname{diag}(z,0,-z)\). The reduction of \(z\) is
\(\bar\Psi(g_1)-\bar\Psi(g_1)^{-1}\ne0\). Lagrange interpolation
gives all three projectors. In rank four the marked block matrix
\(r(\phi)\) has coprime characteristic polynomials on the blocks
\(\{0,2\}\) and \(\{1,3\}\), by the distinct residual values
between blocks. Polynomial interpolation gives the first block
projector. In that block \(r(g_0)+\theta(g_0)r(g_0)^{-1}\) is
diagonal by normalization, with diagonal difference reducing to
\((\bar\alpha(g_0)-1)(\bar\Psi(g_0)-1)\ne0\). We obtain
\(E_{00},E_{22}\) and then their symplectic adjoints \(E_{33},E_{11}\).
\src{src/Development/AutomorphicArithmetic/AutomorphicVarieties/SymplecticGSp4/EiiUnit.lean}{169}{The four marked projectors}.

If the semisimplification is a sum of characters, a full invariant
flag is consequently a coordinate flag: every stable subspace is
preserved by all \(E_{ii}\). Order the coordinates by this flag,
assigning numbers \(b_i\) so \(b_i>b_j\) for every nonzero
off-diagonal entry in position \(ij\). The reverse position
\((j^*,i^*)\), with \(i^*=n-1-i\), also occurs. This follows by
reverse transpose in rank three and symplectic adjunction in rank
four, both preserving \(\mathcal E\). Therefore \(h_i=b_i-b_{i^*}\)
is strictly decreasing on these directed entries. The functional
\(ij\mapsto h_i-h_j\) factors through their weights. For rank three
this follows from \((h_0,h_1,h_2)=(h_0,0,-h_0)\). For rank four it
is the linear form \((h_0-h_1)\lambda+(h_0-h_2)\mu\) on differences
of the coordinates \((\lambda_i,\mu_i)=(1,1),(0,1),(1,0),(0,0)\).
Thus a nonzero balanced monomial uses diagonal entries only. These
are the diagonal characters on the flag and are constant by class
field theory. In rank three all negative-sign elements have the
matrix of a positive-sign element, by multiplication by \(c\).
Conversely, an absolutely irreducible constituent of dimension \(>1\)
forces a nonconstant characteristic-polynomial coefficient if the
residual semisimplification consists of characters. Such coefficients
are themselves sums of weight-zero monomials.

\begin{proposition}[Remaining characteristic-\(\ell\) curves]\label{dp:curvebranches}
For finiteness of the invariant quotient in rank three it suffices
to exclude curves on which \(r|_{G_K}\) is strongly absolutely
irreducible. In rank four it suffices to exclude curves on which
\(r\) is absolutely irreducible, and those with an absolutely
irreducible two-dimensional constituent on which complex conjugation
is scalar.
\end{proposition}
\begin{proof}
The needed local flags specialize even where the marked inertia
eigenvalues coalesce modulo \(\pi\). In rank three take the closed
incidence scheme in the projective full flag variety requiring
stability under \(H\) and the ordered characters \(d_i\) on \(I_0\).
It has a point over the fraction field of the chart constructed
in Proposition~\ref{dp:rank3chart}. Its proper image contains the
generic point of \(R/I_t\) by the kernel assertion there, so every
geometric fibre has a flag. In rank four use the symplectic flag
variety, stability under the full decomposition group and the four
prescribed inertia characters. The condition \(\mathrm{OnComp}\)
gives a domain carrying the local equations and an augmentation to
\(O\) at the diagonal point, through which the test factors. Thus
the differences \(d_i(s_0)-d_j(s_0)\) are nonzero in its fraction
field. There the projectors give a flag by the calculation in
Proposition~\ref{dp:local}. Properness specializes it to the given
geometric point. The closed
conditions in both assertions are intersections of the closed
stability and action equations, so no finite presentation of the
inertia group is needed. Away from \(\ell\) the fixed residual
or finite character kernel bounds inertia. At \(\ell\) the
inertia characters on these flags have finite image in equal
characteristic, also on the full inertia in rank three since
\([I:I_0]\) is finite. These properties pass to subquotients.

In rank three the algebra \(\mathcal E\) contains \(E_{ij}\)
precisely when some matrix entry in that position is nonzero, by
sandwiching between \(E_{ii}\) and \(E_{jj}\). Its entry relation
is thus transitive. Its equivalence classes for mutual
reachability are the simple diagonal blocks after the classes
are topologically ordered, each supporting a full matrix algebra.
Reverse transpose permutes the classes and reverses the order.
The only possible two-element class is \(\{0,2\}\): its reverse
must be the unique class of size two. The middle singleton is
incomparable to that class, since both classes are fixed by
reversal. Thus a two-dimensional constituent forces the actual
outer block decomposition. The middle character and the determinant
on \(G_K\) have finite inertia and trivial reduction and hence are
one by the character argument. The outer block \(U\) has determinant
one. It extends to \(W\) on \(G_{\mathbb Q}\) with
\(W(c)=D=\operatorname{diag}(-1,1)\), by the matrix identity
\(J_2 U(h)^{-t}J_2=D U(h)D^{-1}\) for \(h\in G_K\), with
\(J_2\) the two-by-two antidiagonal unit matrix. The marked
outer coordinates retain the original integral lattice, which
is also preserved by \(D\). The formula on the two open cosets
is continuous, and \(D^2=1\) gives the group law
as in Proposition~\ref{spreal:end}.
Its residual characters are \(\eta\bar\Psi,\bar\Psi^{-1}\).

On \(I_\ell\subset G_K\) the two inertia constituents of \(W\)
are finite and hence constant, equal to \(\bar\Psi,\bar\Psi^{-1}\)
by reduction. They are distinct since \(\psi_\ell^2=1\) and
\(\ell-1\nmid2e\). The restriction to inertia has a common eigenline
by the local flag. The eigenspace for its character is exactly one
line by distinctness. Its character is invariant under the full
decomposition group because it is the restriction of one of the
global residual characters. The line is therefore stable under
that group. Away from \(\ell\) the extension adds only the finite
quadratic character. Lemma~\ref{dp:pan} now excludes an absolutely
irreducible outer block. If an irreducible three-dimensional
representation becomes reducible over a finite number-field
extension, pass to its normal core. Its restriction there is
semisimple: the sum of the simple submodules is nonzero and is
stable under the full group. The isotypic components are permuted
transitively, with equal multiplicities. Primality of dimension
three makes every simple constituent on this subgroup a character.
The preceding potential triangularity argument gives constant
characteristic polynomials, contradicting absolute irreducibility.

In rank four any reducible symplectic similitude representation
has simple constituents of dimension at most two. A stable line
lies in its stable three-dimensional perpendicular and a stable
hyperplane gives such a line by duality. A stable two-plane gives
two pieces of dimension two directly. Every two-dimensional
constituent inherits a local flag by intersection and quotient,
and its residual semisimplification is a pair from
\((\bar\theta,\bar\Psi,\bar\alpha,1)\). Lemma~\ref{dp:pan} excludes
all opposite-parity pairs. The equal-parity pairs are precisely
\(\{\bar\theta,\bar\alpha\}\) and \(\{\bar\Psi,1\}\), with parity
read on the two coordinate pairs \(\{0,2\}\), \(\{1,3\}\).
The preceding marked-entry argument and Lemma~\ref{dp:detect}
complete both reductions.
\end{proof}

The rank-four source separates this issue from its characteristic-zero
purity and automorphy assertions. Its first deduction of finiteness
uses the positive-characteristic exclusion \(\mathrm{Con}(D)\), whose
test includes the nonsplit extension with irreducible even block on
indices \(1,3\). In the tower the ideal \(J\) is required to satisfy
\(\mathfrak m_{X.B}^n\subset J+\mathfrak m_{O'}X.B+
(\psi_{v,i}(s)-1:s\in I_v)\). This gives finiteness on the
fixed-inertia special fibre \emph{of the quotient by \(J\)}.
Using it for a given curve also requires showing that the curve
annihilates \(J\) on the tower. Its localization clause for
irreducible, locally distinguished curves already assumes this
annihilation. The characteristic-zero seed and purity clauses
alone impose no equation on a purely characteristic-\(\ell\)
component. In rank three the source uses a common finite-index
subgroup for constancy of all the special-fibre characteristic
polynomials. The curve criterion shows why it is enough to
settle the remaining representations of Proposition~\ref{dp:curvebranches}
one curve at a time.
\src{src/Development/AutomorphicArithmetic/AutomorphicVarieties/CurveDVRSpecialization/GfResZpow.lean}{179}{Rank-three finiteness from special-fibre characteristic polynomials};
\src{src/Development/AutomorphicArithmetic/AutomorphicVarieties/ModularityArgument/ThresholdArguments/DenseThresholdCases/VWinPm.lean}{75}{The characteristic used in each of the three conclusions};
\src{src/Development/AutomorphicArithmetic/AutomorphicVarieties/ModularityArgument/ThresholdArguments/DenseThresholdCases/PotentialModularityTower.lean}{38}{Finiteness and localization conditions on the tower ideal}.


\subsection{Intersection of the weighted equations}
\label{newcoh:orbitsection}

We prove the geometric inequality used above. Let \(O\) be a complete
discrete valuation ring with uniformizer \(\pi\) and fraction field \(E\)
of characteristic zero. Let \(K\) be an algebraically closed extension
of \(E\) with a real valuation extending the valuation \(v(\pi)=1\).
Write \(A=\bigoplus_{j\in\mathbf Z}A_j\)
and \(B=A_0\). Assume \(B\) is a complete Noetherian local
\(O\)-algebra with the same residue field and local structure map,
and \(A\) is generated over \(B\) by finitely many elements of
weights \(-2,-1,1,2\). An \(O\)-point \(x\) kills all nonzero
weights. Write \(J_x=\ker x\) and use the saturated cotangent module
of Lemma~\ref{newcoh:cotangent}, taking in weight zero the image of
\(J_x\cap B\).

Suppose there is a weight-preserving formal chart
\(A\to E[[X_1,\ldots,X_n]]\) with constant term \(x\), whose linear
terms on \(J_x\) span all coordinate directions over \(E\) and have
kernel the saturation of \(J_x^2\) there. The table gives the number
of chart variables and of homogeneous equations defining an ideal \(I\).
The neutral equations vanish at \(x\).
\[
\begin{array}{c|rrrrr}
 \text{weight}&-1&-2&1&2&0\\ \hline
 \text{variables}&r&s&r'+1&s'+1&m\\
 \text{equations}&r+1&s&r'&s'&m\\
 \text{names}&F&G&P&H&Z
\end{array}
\qquad r\ge1 .
\]
Assume the equations in weights \(-1,-2,0\) span the corresponding
rational cotangent spaces. Set \(X=\operatorname{Spec}(A/I)\).
A geometric point is called two-sided if some positive and some
negative weight both have nonzero value. A positive or negative
one-sided point has nonzero values on precisely one of the two signs.
An invariant value of a point means its restriction to \(B\).

\begin{theorem}[Integral degree of the two-sided orbits]
\label{newcoh:orbitdegree}
Suppose \(B/(I\cap B)\) is finite as an \(O\)-module. Assume every
fixed \(K\)-point of \(X\) other than \(x_K\) has at most one sign of
one-sided points with its invariant value, and the geometric tangent
space of \(X_K\) at each two-sided point has dimension at most one.
Let \(a\in B\) vanish at \(x\) and be nonzero at every two-sided point.
There are finitely many two-sided \(K^\times\)-orbits, and for any finite
set \(T\) of two-sided points meeting every such orbit,
\[
 \operatorname{length}_O Q_{-1}\ \le\ 2\sum_{\xi\in T}v(\xi(a)).
\]
\end{theorem}

\begin{proof}
\emph{The chart and the three gradings.}
Every \(A_j\) is finite over \(B\). Indeed, in a monomial in the
chosen generators one may remove weight-zero submonomials until
none remains. Among exponent vectors of fixed weight there are
only finitely many minimal elements in the componentwise order
by Dickson's lemma, and the removed factors belong to \(B\).
The weight \(-1\) cotangent lattice is thus finite over \(O\)
since \(B\) acts through \(x\) on it. It is torsion-free by
saturation and has rank \(r\) by the chart. Apply Smith normal
form to the map from the \(r+1\) equations to this lattice.
After invertible \(O\)-linear changes of equations and lattice basis
there are \(y_i\in A_{-1}\) and \(c_i\ge0\) such that, modulo
the saturation of \(J_x^2\),
\[
 F_i=\pi^{c_i}y_i\quad(1\le i\le r),\qquad F_*=0,\qquad
 \operatorname{length}_O Q_{-1}=\sum_i c_i .
\]
Here \(F_*\) denotes the last equation. In the formal chart replace
the \(m\) neutral coordinates by the \(Z\)'s using their invertible
linear part and the formal inverse function theorem. Setting those
coordinates to zero gives a chart \(\phi\) with the four nonzero
weights and the same linear coordinates in these weights. The inverse
is weight-preserving by uniqueness of the formal inverse. We may
arrange \(\phi(a)\ne0\). If its image is zero, add \(uF_i\) to \(a\),
choosing \(u\in A_1\) and \(i\le r\) whose images have nonzero linear
parts. The product has nonzero image in a power series domain and
changes no value on \(X\).

Form the following standard \(\mathbf N^3\)-graded subalgebra
of \(A[t_0,t_1,t_2]\), with the ordinary weight in \(A\) kept
separate from the three exponents of \(t\):
\[
 U=A_{-2},\qquad V=A_2,\qquad W=BF_*^2+U^2V,\qquad
 R=B[Ut_0,Vt_1,Wt_2].
\]
Let \(R'\) be its image under the coefficient map \(\phi\), so
\(a\) is a nonzerodivisor on \(R'\). Denote by \(\mathcal S\) the
following list in \(R\) or any of its quotients:
\[
 (F_i^2t_0)_{i\le r},\quad (G_i t_0)_{i\le s},\quad
 (P_i^2t_1)_{i\le r'},\quad (H_i t_1)_{i\le s'},\quad F_*^2t_2 .
\]
Put \(M=R'/(\mathcal S)\). We first determine its horizontal
support, including at the boundary where one sign vanishes.
\src{src/Development/AutomorphicArithmetic/AutomorphicVarieties/CommutativeAlgebra/GradedModulesLength/ToDegInjective.lean}{187}{multigraded algebra and marked equations}.

\emph{Support of high graded pieces.}
Call a prime of \(R\) relevant if it misses some generator in each
of the three standard directions. Let \(\mathfrak p\) be such a prime
above the relations of \(M\), and \(\mathfrak b=\mathfrak p\cap B\).
Choose surviving symbols \(ut_0,vt_1,wt_2\). Since \(W\subset U\),
the identity \((F_*^2t_0)(wt_2)=(wt_0)(F_*^2t_2)\) puts
\(F_*^2t_0\) in \(\mathfrak p\). We have \(I\cap B\subset\mathfrak b\).
Indeed, for a marked coefficient \(h\) of weight \(-2\) and \(k\)
of weight \(2\), multiply \(kh\) by the surviving product
\((ut_0)(vt_1)\). The result is \((ht_0)(kuv t_1)\) in
\(\mathfrak p\). The positive case interchanges the signs. For an
equation of weight \(1\) or \(-1\) apply this argument to its square
and to the square of its complementary coefficient. Neutral equations
are already zero in \(R'\). Taking weight-zero parts of all
multiples of the generators of \(I\) proves the inclusion.

Put \(D=UV\subset B\). We claim that if \(D\subset\mathfrak b\)
then \(\mathfrak b\) is the maximal ideal of \(B\). Otherwise,
since \(B/(I\cap B)\) is finite over \(O\), the point is horizontal
and \(B/\mathfrak b\) embeds in \(K\). We construct one-sided points
of both signs with this invariant value. For \(n\ge1\),
\(A_{2n}=V^n\) and \(A_{-2n}=U^n\). To see this, remove a positive
generator of weight two or a pair of positive generators of weight
one from a monomial of weight \(2n\) until the remaining weight is
two. Negating the weights gives the other equality.
On the even-weight subalgebra define a map to \(\kappa(\mathfrak p)\)
by \(h\mapsto ht_1^n\bmod\mathfrak p\) for \(h\in A_{2n}\), \(n\ge0\),
and send negative weights to zero. This is multiplicative. A mixed
product with nonnegative weight and a negative factor has its
coefficient in \(D V^n\) at weight \(2n\), by the two power
identities. The map kills the appropriate marked even equations and
squares, and takes \(v\) to a nonzero element. The full algebra is
integral over the even-weight subalgebra, because the square of
each odd homogeneous element belongs to it. Lying over extends the
map to a field extension, where all original equations vanish.
Over the invariant value under consideration these equations, the
vanishing of the negative generators and \(v\ne0\) describe a
nonempty finite-type locus over \(\kappa(\mathfrak b)\), hence they
have a \(K\)-solution. The same construction with \(t_0\) gives the
negative point. Taking the scaling limit of either point toward
zero gives the same fixed point. Away from \(x_K\) this contradicts
the hypothesis on the signs. Above \(x_K\) a negative point would
give a curve through \(x_K\), taking \(h\in A_{-j}\) to
\(\xi(h)T^j\), \(j>0\). Choose the least \(j\) occurring with a
nonzero coefficient. Since \(A\) is generated in the stated weights,
\(j=1\) or \(2\). All of \(J_x^2\), and hence its saturation, vanish
to order at least \(2j\) on the curve. The \(F\)'s or \(G\)'s span
the rational cotangent in weight \(-j\) and vanish identically on
the curve, contradicting minimality. This proves the claim.

Here and below \(d\gg0\) means all three components are sufficiently
large. The two prime inclusions just proved give
\[
 (I\cap B)^n M_d=0,\qquad \pi^{n'}M_d\subset D M_d
 \qquad(d\gg0)
\]
for some \(n,n'\). Indeed, multiply the respective ideals \(I\cap B\)
and \((\pi)\) by the irrelevant ideal (the product of the three
direction ideals). The products lie in the radicals of the defining
ideals of \(M\) and \(M/DM\), respectively. Noetherianity gives
powers, and sufficiently high graded pieces are generated by the
corresponding powers of the irrelevant ideal times lower pieces.
The first inclusion makes \(M_d\) finite over \(O\), by the finite
filtration of \(B/(I\cap B)^n\) over \(B/(I\cap B)\). The second
and Nakayama's lemma put its generic support where some
\(u v\ne0\), \(u\in U,\ v\in V\).
\src{src/Development/AutomorphicArithmetic/AutomorphicVarieties/CommutativeAlgebra/GradedModulesLength/SameDegreeZeroOfRestriction.lean}{17}{exclusion of the one-sided boundary}.

\emph{The generic multiplicities.}
After inverting \(uv\), \(u\) is a homogeneous unit in \(A\), and
the three coefficient modules \(U,V,W\) are freely generated over
the localized \(B\) by \(u,u^{-1},u\), respectively. If
\(d=(i,j,k)\) with \(i,j,k\ge1\), taking coefficients identifies
\(R_d\) there with \(A_{2(j-i-k)}\). Multiplication by a marked
section is multiplication by its coefficient between the corresponding
weight spaces. Hence \(M_d\) is a quotient of this weight space
modulo the ideal \(I_2\) obtained from \(I\) by squaring the
\(q=r+1+r'\) equations \(F,P\), and retaining \(G,H,Z\).
Since any product of \(q+1\) generators of \(I\) contains either
a retained equation or a repeated squared generator, \(I^{q+1}\subset I_2\).
In particular all weight spaces occurring in the successive
thickenings are modules over \(B/(I\cap B)^{q+1}\), which is finite
over \(O\) by the same ideal-power filtration. After base change
their local multiplicities may thus be computed over the local
factors of a finite-dimensional \(K\)-algebra.

Before squaring, each weight has geometric multiplicity at most
one at an invariant point where \(uv\ne0\). In fact
\(X_K[(uv)^{-1}]\) is of finite type by finiteness of the invariant
algebra. More precisely, taking the degree-zero summand commutes
with extension of scalars and with localization at the weight-zero
element \(uv\), so its invariant ring is
\(\big((B/(I\cap B))\otimes_O K\big)[1/(uv)]\).
Each closed point of \(X_K[(uv)^{-1}]\) belongs to a one-dimensional
torus orbit, and its tangent space has dimension at most one.
Thus this open scheme is smooth of dimension one there. Its orbits
are open in their components and every closed point lies in such an
orbit, so each connected component is a single orbit. Equivalently,
each orbit is closed here since the invertible coordinate \(u\)
scales by weight \(-2\). The component idempotents are invariant,
so they distinguish the components on the indicated invariant base.
Each weight space on one orbit embeds by pullback into
the corresponding weight space on \(K[t,t^{-1}]\), of dimension at
most one. A finite collection of opens \(uv\ne0\) covers the
two-sided locus, since nonzero weights can be witnessed by
generators and weight-one generators can be squared. This also proves
finiteness of its orbits.

Replacing a homogeneous equation \(f\) of weight \(h\) by \(f^2\)
at most doubles a bound valid for all weights. Indeed, the kernel of
\((L/f^2L)_w\to(L/fL)_w\) is a quotient of \((L/fL)_{w-h}\)
under multiplication by \(f\). Iterate after localizing at the given
geometric invariant point. The multiplicity of \(M_d\otimes_O K\)
at that point is at most \(2^{r+1+r'}\).
On the finite \(O\)-module \(M_d\) put
\(\theta(a,M_d)=\operatorname{length}_O(M_d/aM_d)
-\operatorname{length}_O(M_d[a])\).
Its torsion submodule contributes zero to this difference.
On its free quotient the difference equals the valuation of the
determinant of \(a\), by Smith normal form. Over \(K\) this
determinant is the product of the values of \(a\) at invariant
points counted with the above multiplicities, by decomposition of
the finite-dimensional module into its generalized eigenspaces.
These values are nonzero and integral over \(O\). It follows that
\begin{equation}\label{newcoh:upperorbit}
 \theta(a,M_d)\le 2^{r+1+r'}\sum_{\xi\in T} v(\xi(a)).
\end{equation}
\src{src/Development/AutomorphicArithmetic/AutomorphicVarieties/CommutativeAlgebra/GradedModulesLength/TwoSigned.lean}{17}{multiplicity upper bound for the three gradings}.

\emph{The component through the origin.}
In a chart monomial let \(b_-,b_+\) count its negative and
positive weights with absolute values, so its weight is \(b_+-b_-\).
In degree \(d=(i,j,k)\) of \(R'\) keep the terms with
\[
 (b_-,b_+)=(2i+4k,\,2j+2k).
\]
These are the terms of least possible negative weight. The bounds
for \(U,V\) follow from their weights. For \(W\), the term of
bidegree \((1,0)\) in \(F_*\) is zero since \(F_*\) has zero linear
part, giving \(b_-\ge4\) for its square. The same bound holds on
\(U^2V\). Multiplication of lowest terms therefore defines a
graded quotient \(C\) of \(R'\), retaining the three exponents to
distinguish components. On \(C\) the ring \(B\) acts through
\(x:B\to O\), so \(C\) is a quotient of \(R'/aR'\).
Each \(C_d\) is a finite torsion-free \(O\)-module inside the
polynomials of the indicated bidegree.

Over \(E\) these polynomials all occur. The leading terms in
weight \(-1\) contain all linear coordinates of that weight after
extending scalars. Those from weight \(-2\) span its linear
coordinates modulo quadratics in the weight \(-1\) coordinates.
The quadratics themselves are products of the preceding lifts,
so \(U\) supplies all pure monomials of negative weight two. The
same argument applies to \(V\) on the positive side. A monomial
with even total weight on either side splits into blocks of weight
two, by pairing the weight-one variables. The required blocks form
\(i\) factors from \(U\), \(j\) from \(V\), and \(k\) from
\(U^2V\subset W\). Thus the length of \((C/\pi C)_d\) is
\[
 h(d)=N_{r,s}(i+2k)\,N_{r'+1,s'+1}(j+k),\qquad
 N_{b,c}(z)=[t^{2z}](1-t)^{-b}(1-t^2)^{-c}.
\]
In \(C\) the first \(r\) marked sections are
\(\pi^{2c_i}Y_i\), where \(Y_i\) is the marked leading term of
\(y_i^2\). The Smith relations give this equality because any
weight \(-1\) element in the saturated square has zero linear term.
\src{src/Development/AutomorphicArithmetic/AutomorphicVarieties/CommutativeAlgebra/GradedModulesLength/ReesGradedPieceSurjection.lean}{199}{Hilbert function of the leading terms}.

\emph{Euler characteristics.}
For a finite graded module \(L\) over \(R\), write
\(\chi(\mathbf{s};L)_d\) for the alternating sum of the \(B\)-lengths
of the degree \(d\) Koszul homology for a list of homogeneous
sections. We use this when the high pieces of \(L/(\mathbf{s})L\)
have finite length. Then the high pieces of the homology have finite
length as well, since these finite graded modules are supported on
the same closed locus. Equivalently, off the irrelevant locus all
primes of this support lie over the closed point of \(B\). All
lengths of finite-length \(B\)-modules here agree with their
\(O\)-lengths because the residue fields agree.

For completeness, \(\chi(\mathbf{s};L)_d\ge0\) for \(d\gg0\).
Use induction on the list and a homogeneous prime filtration of \(L\),
using additivity of Koszul Euler characteristics in short exact
sequences. Irrelevant factors vanish in high degree. In a relevant
domain factor, a nonzero first section is regular and the calculation
passes to its quotient with one fewer section. If the first section
vanishes, its two-term complex gives the difference of the Euler
characteristics of the shorter list in degrees \(d\) and
\(d-\deg s_1\). This difference is also the Euler characteristic
of the shorter list on the quotient by any nonzero section of that
degree, by the multiplication exact sequence. Such a section exists
in a relevant domain by taking a product of surviving standard
generators, or \(1\) in degree zero. The shorter list already has
the required support in this case. This proves positivity and,
by additivity, monotonicity on passing to a quotient whenever the
original list has the stated support.

The support condition holds for \((\mathcal S,a)\) on \(R'\) by
the preceding description of \(M_d\) and the nonvanishing of \(a\)
on its generic support. Koszul complexes are invariant up to
isomorphism under permutation. Since \(a\) is regular on \(R'\),
monotonicity and the one-section kernel--cokernel formula give
\[
 \chi(\mathcal S;C)_d
 \le\chi(\mathcal S,a;R')_d
 \le\chi(a;M)_d=\theta(a,M_d)\qquad(d\gg0).
\]
For the second inequality successively replace each of the first
sections by its ordinary quotient. At each step the omitted term
is the Euler characteristic of the remaining list on the shifted
kernel of that multiplication and is nonnegative by the preceding
argument.

In \(C\), split each \(\pi^{2c_i}Y_i\) into its two factors using
additivity of the one-section Koszul class for a product. In fact,
writing \(e_0=(1,0,0)\), use the degree-preserving composite
\[
 C(-e_0)\xrightarrow{\,Y_i\,}C\xrightarrow{\,\pi^{2c_i}\,} C .
\]
The three mapping cones for the first map, the composite, and the
second map form an exact triangle in that order. Tensor with the
remaining Koszul complex, whose terms are graded free, and take
degree \(d\). Taking Euler characteristics expresses the contribution
of the composite as the sum of those of the two factors, with
both summands taken in degree \(d\). The same triangle splits the
degree-zero power into \(2c_i\) copies of the section \(\pi\).
Each factor list has support contained in
the original intersection. In each term with a single \(\pi\),
regularity of \(\pi\) on \(C\) passes to \(C/\pi C\).
Its terms already have finite length degree by degree, so the Euler
characteristic is simply the alternating sum of the lengths of its
terms, independent of the actual sections of the remaining degrees.
Peeling off the powers successively leaves one final nonnegative
Euler characteristic. With \(\Delta_i,\Delta_j,\Delta_k\) denoting
unit backward differences in the indicated variable this gives
\[
 \chi(\mathcal S;C)_d\ge
 2\left(\sum_i c_i\right)
 \Delta_i^{\,r+s-1}\Delta_j^{\,r'+s'}\Delta_k h(d).
\]
For \(b\ge1\), partial fractions at \(t=1,-1\) show that
\(N_{b,c}(z)\) is eventually polynomial of degree \(b+c-1\)
with leading coefficient \(2^{b-1}/(b+c-1)!\). Indeed the pole
at \(1\) has order \(b+c\) and leading factor \(2^{-c}\);
the pole at \(-1\) has order at most \(c\), of lower degree.
The indicated \(i\)-difference therefore gives \(2^{r-1}\) on
the first factor, and the \(j\)-difference gives a linear polynomial
of slope \(2^{r'}\) in \(j+k\) on the second. Its \(k\)-difference
is \(2^{r'}\). We have proved
\[
 2^{r+r'}\sum_i c_i\le \theta(a,M_d)
 \le 2^{r+1+r'}\sum_{\xi\in T}v(\xi(a))
 \qquad(d\gg0).
\]
Dividing gives the theorem.
\end{proof}
\src{src/Development/AutomorphicArithmetic/AutomorphicVarieties/CommutativeAlgebra/GradedModulesLength/AlgMemCoef.lean}{92}{comparison of the two Koszul calculations}.

For the arithmetic application the valuation on a finite coefficient
extension has \(v(z)=-e(O/\mathbf Z_\ell)\log|z|_\ell/\log\ell\).
Thus Lemma~\ref{newcoh:cotangent} and Theorem~\ref{newcoh:orbitdegree}
give the factor \(2\) in the normalized trace inequality once the
rank-four deformation datum supplies the chart, equations, one-sided
exclusion, infinitesimal rigidity and trace function with the stated
nonvanishing. Establishing these hypotheses for all cases II and III
and bounding the resulting automorphic family remain the global steps.
 
